\chapter{Differential and Integral Calculus}\label{ch:calculus}

\begin{center}\itshape
$\dfrac\pi2=\dfrac21\cdot\dfrac23\cdot\dfrac43\cdot\dfrac45\cdot\dfrac65\cdot\dfrac67\cdots$ \quad and \quad $\displaystyle\int_0^\infty e^{-x^2}\dd x=\frac{\sqrt\pi}2$.\\
Both come out of one humble family of integrals, which also completes Stirling's formula.
\end{center}

In 1656 John Wallis, Savilian professor at Oxford, wanted a value for $\pi$. He computed $\int_0^1(1-x^2)^n\,\dd x$ for whole numbers $n$. Each value is a ratio of factorials. He guessed how the formula should continue at half-integers, where the integral is the area of a quarter disc, and wrote an infinite product for $\pi/2$. He had no proof. The \emph{Arithmetica Infinitorum} says so. The guess irritated Fermat, who thought induction from a few cases was no way to do mathematics, and the two quarrelled by letter for a year.

Here the same family of integrals, now named after Wallis, is proved in a page. The same recurrence then delivers the constant $\sqrt{2\pi}$ in Stirling's formula and the Gaussian integral. Three famous results, one humble relation.

The other thread is Taylor's formula. It replaces a function by a polynomial and says exactly how much is lost. Brook Taylor published the expansion in 1715, without an error estimate and without noticing that it contained the series for $\sin$, $\cos$ and $\exp$. Lagrange later called it the principal foundation of the differential calculus and, in 1797, gave the remainder a usable form. The Gaussian $e^{-x^2}$ enters as de Moivre's 1733 approximation to binomial coefficients. Laplace evaluated its integral in 1774.

The order of work is forced by those stories. Derivatives and the fundamental theorem come first, because every later estimate is an integral of a derivative. Wallis's integrals are the worked example. L'Hôpital's rule and the Gamma function are two applications. Numerical integration is the place where the remainder stops being a symbol and becomes a number of decimals. Along the way, Euler's Gamma function interpolates the factorials, as in his letter to Goldbach of 1729, and two short integrations by parts prove that $e$ and $\pi$ are irrational: Fourier's argument of about 1815, and Niven's of 1947.

\portraits{taylor}{Brook Taylor (1685--1731). Portrait, artist unknown. Public domain, via Wikimedia Commons.}{wallis}{John Wallis (1616--1703). After Godfrey Kneller, 1701. Public domain, via Wikimedia Commons.}

\section{Derivatives and antiderivatives, from scratch}

We write $f^{(k)}$ for the $k$th derivative, with $f^{(0)}=f$.
A function is of class $C^k$ on an open interval if its derivatives through
order $k$ exist and are continuous there; $C^0$ means continuous.
On a closed interval the derivatives are continuous up to the endpoints,
with endpoint derivatives interpreted one-sided. These hypotheses make the
suprema in the remainder estimates finite.

Chapter~\ref{ch:sequences} defined the derivative and proved the mean value theorem. The rules for computing derivatives are the next thing a working mathematician
needs, and they are quick to prove.

\begin{propositionx}[rules of differentiation]
Let $f$ and $g$ be differentiable at an interior point $a$ of their intervals of definition.
\begin{enumerate}[label=(\alph*)]
\item $(f+g)'(a)=f'(a)+g'(a)$ and $(fg)'(a)=f'(a)g(a)+f(a)g'(a)$; if $g(a)\ne0$, $\left(\frac fg\right)'(a)=\frac{f'(a)g(a)-f(a)g'(a)}{g(a)^2}$.
\item \textbf{Chain rule}: if $h$ is differentiable at $f(a)$, then $(h\circ f)'(a)=h'(f(a))\,f'(a)$.
\item \textbf{Inverse functions}: if $f$ is continuous and strictly monotone on an interval $I$, with $a$ interior to $I$ and $f'(a)\ne0$, then $f^{-1}$ is differentiable at $b=f(a)$ with $(f^{-1})'(b)=\frac1{f'(a)}$.
\end{enumerate}
\end{propositionx}
\begin{proof}
(a)--(b) These are Proposition~\ref{prop:early-differentiation}, proved in
Chapter~\ref{ch:sequences} before its first applications to series and integrals.

(c) $f$ is a bijection from $I$ onto the interval $J=f(I)$ (Theorem~\ref{thm:ivt}), and $g=f^{-1}:J\to I$ is monotone. A monotone function whose image is an interval is continuous: if, say, $g$ is increasing and discontinuous at $b$, its one-sided limits $g(b^-)\le g(b)\le g(b^+)$ exist (they are suprema and infima of bounded monotone pieces) and are not all equal, so the values strictly between $g(b^-)$ and $g(b^+)$ other than $g(b)$ are missed by $g$, and the image $I$ is not an interval. Then
$\frac{f^{-1}(y)-f^{-1}(b)}{y-b}=\frac{x-a}{f(x)-f(a)}$ with $x=f^{-1}(y)\to a$ as $y\to b$, and the quotient tends to $\frac1{f'(a)}$.
\end{proof}

\begin{reminder}[what is admitted about the elementary functions]
The volume takes for granted the functions $\exp$, $\ln$, $\cos$, $\sin$ and the number $\pi$ with their basic properties: $\exp'=\exp$, $\exp(0)=1$, $\exp(a+b)=\exp a\exp b$, $\exp$ is a bijection from $\R$ onto $]0,\infty[$ with inverse $\ln$, so $\ln'(x)=\frac1x$ by (c); $\sin'=\cos$, $\cos'=-\sin$, $\cos^2+\sin^2=1$, the addition formulas, and $\pi$ is twice the first positive zero of $\cos$. The series $\sum\frac{x^n}{n!}$ is one route to constructing the exponential. Its absolute convergence follows from the ratio test; Exercise~\ref{ex:taylorln}(c)
identifies it with the complex exponential defined in Chapter~\ref{ch:complex}. We use the elementary-function properties listed here as named assumptions. From these, the power functions $x^\alpha=e^{\alpha\ln x}$ ($x>0$) have derivative $\alpha x^{\alpha-1}$, and $\tan=\frac{\sin}{\cos}$ has derivative $1+\tan^2$; its inverse $\arctan:\R\to\,]-\frac\pi2,\frac\pi2[$ has derivative $\frac1{1+x^2}$ by (c).
\end{reminder}

\begin{examplex}[three ways to find an antiderivative]
An \emph{antiderivative} of $f$ on an interval is a function $F$ with $F'=f$; two antiderivatives differ by a constant (Theorem~\ref{thm:mvt}(c)), and
Proposition~\ref{prop:inttools} below shows that every continuous function has one, namely $x\mapsto\int_a^xf$. Three techniques find them in practice; each is a
differentiation rule read backwards.

\emph{Substitution} (the chain rule backwards): to find $\int\frac{x}{1+x^2}\dd x$, notice $x\dd x=\frac12\dd(1+x^2)$, so the integral is $\frac12\ln(1+x^2)+C$. To find $\int_0^1\sqrt{1-x^2}\dd x$, put $x=\sin\theta$, $\dd x=\cos\theta\dd\theta$: the integral becomes $\int_0^{\pi/2}\cos^2\theta\dd\theta=\frac\pi4$ (the area of a quarter disc, as it should be).

\emph{Integration by parts} (the product rule backwards): to find $\int x\,e^x\dd x$, differentiate the polynomial and integrate the exponential: $xe^x-\int e^x\dd x=(x-1)e^x+C$. To find $\int\ln x\dd x$, write it as $\int1\cdot\ln x\dd x$ and differentiate the logarithm: $x\ln x-\int x\cdot\frac1x\dd x=x\ln x-x+C$. The rule of thumb: differentiate what gets simpler (polynomials, logarithms, inverse trigonometric functions), integrate what stays the same (exponentials, sines, cosines).

\emph{Partial fractions} (for rational functions): to find $\int\frac{\dd x}{x^2-1}$, write $\frac1{x^2-1}=\frac12\left(\frac1{x-1}-\frac1{x+1}\right)$, so the integral is $\frac12\ln\left|\frac{x-1}{x+1}\right|+C$. Quadratic denominators may instead lead to $\arctan$: completing the square reduces $\int\frac{\dd x}{x^2+px+q}$, when $4q>p^2$, to a rescaled $\int\frac{\dd u}{1+u^2}$. Each proposed antiderivative can be checked by differentiation.
\end{examplex}

\begin{reminder}[improper integrals]
When $f$ is continuous on $[a,\infty[$, we define $\int_a^\infty f=\lim_{A\to\infty}\int_a^Af$ when the limit exists, and say the integral \emph{converges}. Likewise $\int_a^bf=\lim_{\varepsilon\to0^+}\int_{a+\varepsilon}^bf$ when $f$ is continuous on $]a,b]$ only. For $f\ge0$ the integral converges if and only if the partial integrals are bounded (they are non-decreasing in $A$), and comparison works as for series: $0\le f\le g$ and $\int g<\infty$ give $\int f<\infty$. Thus $\int_1^\infty\frac{\dd x}{x^\alpha}$ converges iff $\alpha>1$, $\int_0^1\frac{\dd x}{x^\alpha}$ iff $\alpha<1$, $\int_0^\infty e^{-x}\dd x=1$, and $\int_0^\infty e^{-x^2}\dd x$ converges because $e^{-x^2}\le e^{-x}$ for $x\ge1$. If $\int|f|$ converges, so does $\int f$: for large $A<B$,
$|\int_A^B f|\le\int_A^B|f|$ is arbitrarily small. Thus truncated integrals
along integer endpoints form a Cauchy sequence, and the same estimate shows
that every real endpoint tends to its limit. The argument near a finite
improper endpoint is identical after substituting a sequence of cutoffs. For example, on $[1,A]$ integration by parts gives
$\int_1^A\sin x/x\,\dd x=[-\cos x/x]_1^A-\int_1^A\cos x/x^2\,\dd x$.
The boundary term converges and the last integral converges absolutely by comparison
with $1/x^2$. At zero, $\sin x/x$ extends continuously by $1$.
Absolute convergence fails: on $[k\pi,(k+1)\pi]$, $k\ge1$, the integral of
$|\sin x|/x$ is at least $2/((k+1)\pi)$, whose sum diverges.
Thus $\int_0^\infty\sin x/x\,\dd x$ is a concrete conditional example;
its value $\pi/2$ is an advanced outlook proved in Volume~4.
\end{reminder}

\section{Tools}

We retain the integration contract of Chapter~\ref{ch:sequences}: continuous
functions on closed bounded intervals have an integral with linearity,
monotonicity, additivity and the stated values for constants.
Proposition~\ref{prop:early-ftc} proved the fundamental theorem from these
assumptions; we recall it below to collect the working tools in one place. (Volume~6 constructs the integral; the construction uses Heine's theorem, which Chapter~\ref{ch:sequences} proved.)

\begin{proposition}[integration toolkit]\label{prop:inttools}
Let $f$ be continuous on an interval $I$ and $a\in I$.
\begin{enumerate}[label=(\alph*)]
\item \textbf{Fundamental theorem}: $F(x)=\int_a^xf$ is differentiable with $F'=f$ at every interior point of $I$. Hence $\int_a^bg'=g(b)-g(a)$ for $g$ of class $C^1$.
\item \textbf{Integration by parts}: for $u,v$ of class $C^1$ on $[a,b]$, $\int_a^bu'v=[uv]_a^b-\int_a^buv'$.
\item \textbf{Change of variables}: for $\varphi$ of class $C^1$ on $[a,b]$ with values in $I$,
\[\int_{\varphi(a)}^{\varphi(b)}f(x)\dd x=\int_a^bf(\varphi(t))\,\varphi'(t)\dd t.\]
\item \textbf{Riemann sums}: for $f$ of class $C^1$ on $[0,1]$,
\[\left|\frac1n\sum_{k=1}^nf\left(\frac kn\right)-\int_0^1f\right|\le\frac{\sup|f'|}{n}.\]
\end{enumerate}
\end{proposition}
\begin{proof}
(a) This is Proposition~\ref{prop:early-ftc}; its proof used only the
integration contract, continuity and the mean value theorem.

(b) Apply (a) to $(uv)'=u'v+uv'$, continuous.

(c) Both sides, as functions of the upper limit $b$, vanish at $b=a$ and have the same derivative $f(\varphi(b))\varphi'(b)$, by (a) and the chain rule.

(d) On $\left[\frac{k-1}n,\frac kn\right]$, $|f(t)-f(\frac kn)|\le\frac{\sup|f'|}n$ by the mean value theorem, so
\[\left|\int_{(k-1)/n}^{k/n}f-\frac1nf\left(\frac kn\right)\right|\le\frac{\sup|f'|}{n^2}.\]
Sum over $k$ (Figure~\ref{fig:riemannsum}, left).
\end{proof}

Statement (d) holds for every continuous $f$, with an error that tends to $0$ but is no longer quantified: by Heine's theorem (Theorem~\ref{thm:heine}), given $\varepsilon>0$ there is $\delta$ with $|f(t)-f(s)|\le\varepsilon$ whenever $|t-s|\le\delta$, and for $n>\frac1\delta$ the same computation gives an error at most $\varepsilon$. This is the whole reason Heine's theorem was proved: it is what makes the integral of a continuous function a limit of finite sums.

\begin{figure}[ht]
\centering
\fig{ch06-riemann}
\caption{Left: the Riemann sum of $f(x)=e^{-x^2}$ with $n=8$ right endpoints (bars) against the area under the curve. The error, the sum of the small curved triangles, is at most $\frac{\sup|f'|}n\approx0.11$ by (d); the actual error is $0.040$. Right: the functions $\sin^nt$ on $[0,\frac\pi2]$ for $n=1,2,5,20,100$. The graph collapses towards $t=\frac\pi2$, and the width of the bump shrinks like $\frac1{\sqrt n}$, which is why $W_n\sim\sqrt{\pi/2n}$ (Proposition~\ref{prop:wallis}).}
\label{fig:riemannsum}
\end{figure}

\begin{think}
Without a calculator, estimate $\sin(0.1)$. The tangent at $0$ gives $0.1$; how far off is it, and what would you add to do better? How many terms of the form $\pm\frac{x^k}{k!}$ do you need for six correct decimals?
\end{think}

\begin{theorem}[Taylor formulas]\label{thm:taylor}
Let $f$ be of class $C^{n+1}$ on an interval containing $a$ and $x$, and $T_n(x)=\sum_{k=0}^n\frac{f^{(k)}(a)}{k!}(x-a)^k$.
\begin{enumerate}[label=(\alph*)]
\item \textbf{Integral remainder}:
\[f(x)=T_n(x)+\int_a^x\frac{(x-t)^n}{n!}f^{(n+1)}(t)\dd t.\]
\item \textbf{Taylor--Lagrange inequality}:
\[|f(x)-T_n(x)|\le\frac{|x-a|^{n+1}}{(n+1)!}\sup|f^{(n+1)}|,\]
the supremum over $J=[\min(a,x),\max(a,x)]$.
\item \textbf{Taylor--Young}: if $f$ is only $C^n$, $f(x)=T_n(x)+o\left((x-a)^n\right)$ as $x\to a$.
\end{enumerate}
\end{theorem}
\begin{proof}
(a) Induction on $n$. For $n=0$ it is Proposition~\ref{prop:inttools}(a). If it holds at order $n-1$ for a $C^{n+1}$ function, integrate the remainder by parts with
\[
\begin{gathered}u=-\frac{(x-t)^n}{n!},\qq u'=\frac{(x-t)^{n-1}}{(n-1)!},\qq v=f^{(n)}:\end{gathered}
\]
\[
\begin{aligned}&\int_a^x\frac{(x-t)^{n-1}}{(n-1)!}f^{(n)}(t)\dd t\eqbr\quad=\frac{(x-a)^n}{n!}f^{(n)}(a)+\int_a^x\frac{(x-t)^n}{n!}f^{(n+1)}(t)\dd t.\end{aligned}
\]
(b) Let $M=\sup_{t\in J}|f^{(n+1)}(t)|$, with $J=[\min(a,x),\max(a,x)]$. The integrand of (a) is at most $M\frac{|x-t|^n}{n!}$ in absolute value, and
\[
\left|\int_a^x\frac{|x-t|^n}{n!}\dd t\right|=\frac{|x-a|^{n+1}}{(n+1)!},
\]
so $|f(x)-T_n(x)|\le M\frac{|x-a|^{n+1}}{(n+1)!}$.

(c) For $n=0$ the claim is continuity. For $n\ge1$, apply (a) at order $n-1$ and subtract
\[\frac{f^{(n)}(a)}{n!}(x-a)^n=\int_a^x\frac{(x-t)^{n-1}}{(n-1)!}f^{(n)}(a)\dd t.\]
This gives
\[\begin{aligned}|f(x)-T_n(x)|&\le\frac{|x-a|^n}{n!}\eqbr\times\sup_{t\in[\min(a,x),\max(a,x)]}\left|f^{(n)}(t)-f^{(n)}(a)\right|,\end{aligned}\]
and the sup tends to $0$ by continuity of $f^{(n)}$.
\end{proof}

\begin{examplex}[$\sin(0.1)$, as promised]
For $f=\sin$ at $a=0$, every derivative is $\pm\sin$ or $\pm\cos$, bounded by $1$, so Taylor--Lagrange gives $|\sin x-T_n(x)|\le\frac{|x|^{n+1}}{(n+1)!}$. The tangent $T_1(x)=x$ gives
$\sin(0.1)\approx0.1$ with error at most $\frac{0.01}2=0.005$ (in fact $0.000167$: the bound is crude because $T_1=T_2$, so one may use $n=2$ and get $\frac{0.001}6$). With $T_3(x)=x-\frac{x^3}6$: $0.1-0.000166\overline6=0.09983\overline3$, error at most $\frac{10^{-5}}{120}<10^{-7}$ (again $n=4$ works): six correct decimals already, $\sin(0.1)=0.0998334\ldots$ Two terms. The Taylor series converges to $\sin x$ for every $x$ (Figure~\ref{fig:taylor}, left), but for $x=\pi$ one needs $n$ well past $\pi$ before the bound $\frac{\pi^{n+1}}{(n+1)!}$ becomes small; a low-degree Taylor polynomial is most useful near its centre; increasing the degree enlarges the region of good approximation for sine.
\end{examplex}

Figure~\ref{fig:taylor} (left) shows the Taylor polynomials of $\sin$ at $0$ hugging the curve on a wider and wider window. Taylor--Lagrange explains why: the bound $\frac{|x|^{n+1}}{(n+1)!}$ tends to $0$ for each fixed $x$, but only once $n$ is well past $|x|$. Figure~\ref{fig:taylorexp} shows two other fates: for $\exp$ the polynomials converge everywhere, for $\ln(1+x)$ they converge on $]-1,1]$ and nowhere else, however many terms one takes.

\begin{figure}[ht]
\centering
\fig{ch06-taylor}
\caption{Left: $\sin$ and its Taylor polynomials $T_1,T_5,T_9,T_{13}$ at $0$. At $x=\pi$ they give $3.14$, $0.524$, $0.0069$ and $0.00002$, closing in on $\sin\pi=0$. Right: the sandwich of Exercise~\ref{ex:gaussian} for $n=3$. The three areas are $\sqrt3\,W_7\approx0.792$, $\frac{\sqrt\pi}2\approx0.886$ and $\sqrt3\,W_4\approx1.020$. As $n$ grows, both bounding areas close in on the Gaussian integral $\sqrt\pi/2$.}
\label{fig:taylor}
\end{figure}

\begin{figure}[ht]
\centering
\fig{ch06-taylorexp}
\caption{Left: $e^x$ and its Taylor polynomials $T_1,\dots,T_4$ at $0$; for each fixed $x$ these converge to $e^x$. Successive degrees need not improve the error monotonically when $x<0$. Right: $\ln(1+x)$ and its Taylor polynomials $T_2,T_5,T_{10},T_{20}$ at $0$. On $]-1,1[$ they converge; at $x=1$ they still do, slowly (Exercise~\ref{ex:taylorln}); for fixed $x>1$ the series diverges because its terms do not tend to zero. A particular partial sum may nevertheless happen to approximate the value. The series $\sum\frac{(-1)^{k+1}x^k}k$ has radius of convergence $1$, a notion Volume~4 studies.}
\label{fig:taylorexp}
\end{figure}

\section{Wallis integrals}

For $n\in\N$ let $W_n=\displaystyle\int_0^{\pi/2}\sin^nt\dd t$ (equal to $\int_0^{\pi/2}\cos^nt\dd t$ by $t\mapsto\frac\pi2-t$).

\begin{think}
Sketch $\sin^nt$ on $\left[0,\frac\pi2\right]$ for $n=1,10,100$. The graph collapses towards $\frac\pi2$, so $W_n\to0$. How fast? Guess the width of the region where $\sin^nt$ is not small, using $\sin\left(\frac\pi2-s\right)=\cos s\approx e^{-s^2/2}$. Figure~\ref{fig:riemannsum} (right) shows the answer.
\end{think}

\begin{proposition}\label{prop:wallis}
\begin{enumerate}[label=(\roman*)]
\item $W_0=\frac\pi2$, $W_1=1$, and $W_n=\frac{n-1}nW_{n-2}$ for $n\ge2$.
\item $W_{2p}=\dfrac{(2p)!}{4^p(p!)^2}\cdot\dfrac\pi2$ and $W_{2p+1}=\dfrac{4^p(p!)^2}{(2p+1)!}$.
\item $(W_n)$ is decreasing and positive, $nW_nW_{n-1}=\frac\pi2$, and $W_n\sim\sqrt{\dfrac\pi{2n}}$.
\end{enumerate}
\end{proposition}
\begin{proof}
(i) Integrate by parts with $u'=\sin t$, $v=\sin^{n-1}t$, so $u=-\cos t$ and the boundary term $[-\cos t\sin^{n-1}t]_0^{\pi/2}$ vanishes:
$W_n=(n-1)\int_0^{\pi/2}\cos^2t\sin^{n-2}t\dd t=(n-1)(W_{n-2}-W_n)$, using $\cos^2=1-\sin^2$.

(ii) Iterate (i): $W_{2p}=\frac{(2p-1)(2p-3)\cdots1}{(2p)(2p-2)\cdots2}W_0$, and multiply numerator and denominator by $(2p)(2p-2)\cdots2=2^pp!$. Same for odd indices.

(iii) $0\le\sin^{n}t\le\sin^{n-1}t$ on $\left[0,\frac\pi2\right]$, so $(W_n)$ decreases, and $W_n>0$ since the integrand is positive on $]0,\frac\pi2]$. By (i), $nW_nW_{n-1}=(n-1)W_{n-1}W_{n-2}$, so this quantity is constant, equal to $1\cdot W_1W_0=\frac\pi2$.
Finally $\frac{n-1}n=\frac{W_n}{W_{n-2}}\le\frac{W_n}{W_{n-1}}\le1$, so $W_n\sim W_{n-1}$ and $W_n^2\sim nW_nW_{n-1}/n=\frac\pi{2n}$.
\end{proof}

\begin{think}
Compute $\frac{10!}{\sqrt{10}\,(10/e)^{10}}$ on a calculator. You should get $2.5276$. Which famous constant is close?
\end{think}

\begin{theorem}[Stirling's constant]\label{thm:stirlingC}
The constant of Exercise~\ref{ex:stirling} is $C=\sqrt{2\pi}$, so $n!\sim\sqrt{2\pi n}\left(\frac ne\right)^n$.
\end{theorem}
\begin{proof}
Insert $m!\sim C\sqrt m\left(\frac me\right)^m$ into (ii):
\[
\begin{aligned}W_{2p}&=\frac\pi2\cdot\frac{(2p)!}{4^p(p!)^2}\eqbr\sim\frac\pi2\cdot\frac{C\sqrt{2p}\,(2p)^{2p}e^{-2p}}{4^pC^2p\,p^{2p}e^{-2p}}\eqbr=\frac\pi2\cdot\frac{\sqrt2}{C\sqrt p}.\end{aligned}
\]
By (iii), $W_{2p}\sim\sqrt{\frac\pi{4p}}$. Comparing, $\frac{\pi\sqrt2}{2C}=\frac{\sqrt\pi}2$, so $C=\sqrt{2\pi}$.
\end{proof}

De Moivre's $2.5074$ was $\sqrt{2\pi}=2.5066$ to three decimals; the box above, with $n=10$, gives $2.5276$, which is $\sqrt{2\pi}\left(1+\frac1{120}\right)=2.5275$ within $10^{-4}$. That $\frac1{12n}$ is no accident: Proposition~\ref{prop:stirlingnext} below explains it.

\section{L'H\^opital's rule}

The first calculus textbook ever printed, the \emph{Analyse des infiniment petits} of 1696, carries the name of the Marquis de l'H\^opital, and so does
its most famous rule; but the marquis had bought it. In 1694 he had agreed to pay Johann Bernoulli a yearly pension of 300 livres in exchange for
Bernoulli's discoveries and the right to use them as he liked, and the rule for $\frac00$ was among the goods delivered. Bernoulli complained after
l'H\^opital's death; the contract, found in 1921, proved him right. The rule itself is a corollary of a generalised mean value theorem due to Cauchy.

\begin{theorem}[l'H\^opital's rule]\label{thm:lhopital}
\begin{enumerate}[label=(\alph*)]
\item Let $a<b$ and let $f,g$ be continuous on $[a,b]$ and differentiable on $]a,b[$. Then, for some $c\in\,]a,b[$,
\[\begin{aligned}\bigl(f(b)-f(a)\bigr)g'(c)\eqbr=\bigl(g(b)-g(a)\bigr)f'(c).\end{aligned}\]
\item Let $f,g$ be differentiable on $]a,a+h[$ with $g'\ne0$ there, $f(x)\to0$ and $g(x)\to0$ as $x\to a^+$. If $\frac{f'(x)}{g'(x)}\to\ell$ as $x\to a^+$ (with $\ell\in\R$ or $\ell=\pm\infty$), then $\frac{f(x)}{g(x)}\to\ell$ as well. The same holds for $x\to a^-$, for two-sided limits, and for $x\to\pm\infty$.
\end{enumerate}
\end{theorem}
\begin{proof}
(a) Apply Rolle (Theorem~\ref{thm:mvt}) to $\varphi(x)=\bigl(f(b)-f(a)\bigr)g(x)-\bigl(g(b)-g(a)\bigr)f(x)$, which is continuous on $[a,b]$, differentiable inside, and takes the value $f(b)g(a)-g(b)f(a)$ at both ends.

(b) Extend $f$ and $g$ by $f(a)=g(a)=0$: they become continuous on $[a,a+h[$. Note $g(x)\ne0$ for $x\in\,]a,a+h[$: otherwise Rolle would give a zero of $g'$. For such $x$, (a) on $[a,x]$ gives $c_x\in\,]a,x[$ with $f(x)g'(c_x)=g(x)f'(c_x)$, that is $\frac{f(x)}{g(x)}=\frac{f'(c_x)}{g'(c_x)}$. As $x\to a^+$, $c_x\to a^+$ too, so the right side tends to $\ell$. The left-sided and two-sided cases are the same; for $x\to+\infty$ substitute $x=\frac1t$ and apply the result at $t\to0^+$ (the chain rule turns $f'(1/t)$ into $-\frac1{t^2}f'(1/t)$ in both numerator and denominator, and the factors cancel).
\end{proof}

\begin{examplex}
$\frac{\sin x}x\to1$ as $x\to0$: the derivatives give $\frac{\cos x}1\to1$. $\frac{1-\cos x}{x^2}\to\frac12$: differentiate once to get $\frac{\sin x}{2x}$, then once more to get $\frac{\cos x}2$ (the rule may be applied repeatedly as long as the hypotheses hold). $\frac{\ln x}{x}\to0$ as $x\to\infty$: the quotient of derivatives is $\frac1x\to0$; this is the ``$\frac\infty\infty$'' form, which has its own version of the rule, with the same proof (Exercise~\ref{ex:lhopital}). Beware: when $\frac{f'}{g'}$ has no limit the rule says nothing, and $\frac fg$ may still converge: $\frac{x+\sin x}x\to1$ as $x\to\infty$, but $\frac{1+\cos x}1$ oscillates. And Taylor--Young is usually faster: $\frac{e^x-1-x}{x^2}=\frac{\frac{x^2}2+o(x^2)}{x^2}\to\frac12$ in one line.
\end{examplex}

\section{The Gamma function}

On 13 October 1729 Euler wrote to Christian Goldbach that he had found a formula ``interpolating'' the factorials: a function of a
continuous variable $x$ that equals $n!$ at $x=n$. The letter gives an infinite product; by the next year he had the integral below, which is the form everybody uses.

\begin{proposition}[the Gamma function]\label{prop:gammafn}
For $x>0$ the integral $\Gamma(x)=\displaystyle\int_0^\infty t^{x-1}e^{-t}\dd t$ converges, and
\begin{enumerate}[label=(\alph*)]
\item $\Gamma(x+1)=x\,\Gamma(x)$ for $x>0$; hence $\Gamma(n+1)=n!$ for $n\in\N$.
\item $\Gamma\left(\frac12\right)=\sqrt\pi$, and more generally $\Gamma\left(n+\frac12\right)=\dfrac{(2n)!}{4^nn!}\sqrt\pi$.

\end{enumerate}
\end{proposition}
\begin{proof}
Convergence. Near $0$: $t^{x-1}e^{-t}\le t^{x-1}$, and $\int_0^1t^{x-1}\dd t=\frac1x$ converges since $x-1>-1$. Near $\infty$: $t^{x-1}e^{-t}\le e^{-t/2}$ for $t$ large (because $t^{x-1}e^{-t/2}\to0$, as $\frac{\ln t}{t}\to0$), and $\int_1^\infty e^{-t/2}\dd t$ converges.

(a) Integrate by parts on $[\varepsilon,A]$, differentiating $t^x$ and integrating $e^{-t}$: $\int_\varepsilon^At^xe^{-t}\dd t=\left[-t^xe^{-t}\right]_\varepsilon^A+x\int_\varepsilon^At^{x-1}e^{-t}\dd t$. As $\varepsilon\to0$ and $A\to\infty$ the bracket tends to $0$ (for $x>0$), and the two integrals tend to $\Gamma(x+1)$ and $\Gamma(x)$. Then $\Gamma(1)=\int_0^\infty e^{-t}\dd t=1$ and induction give $\Gamma(n+1)=n!$.

(b) Substitute $t=u^2$, $\dd t=2u\dd u$: $\Gamma\left(\frac12\right)=\int_0^\infty t^{-1/2}e^{-t}\dd t=2\int_0^\infty e^{-u^2}\dd u=\sqrt\pi$ by Exercise~\ref{ex:gaussian}. Then (a) repeatedly: $\Gamma\left(n+\frac12\right)=\left(n-\frac12\right)\left(n-\frac32\right)\cdots\frac12\,\Gamma\left(\frac12\right)=\frac{(2n-1)(2n-3)\cdots1}{2^n}\sqrt\pi$, and $(2n-1)(2n-3)\cdots1=\frac{(2n)!}{2^nn!}$.

\end{proof}

So $\left(\frac12\right)!$, if it means anything, is $\Gamma\left(\frac32\right)=\frac{\sqrt\pi}2\approx0.886$, which is exactly Wallis's interpolated value of 1656. Figure~\ref{fig:gammafn}
shows the graph: a minimum near $x=1.46$, the factorials on the integers, and a pole at $0$. Chapter~\ref{ch:inequalities}, Proposition~\ref{prop:gamma-convexity}, proves
that $\ln\Gamma$ is convex. As an advanced outlook, Bohr and Mollerup
proved in 1922 that $\Gamma$ is the \emph{only} function on $]0,\infty[$ with $\Gamma(1)=1$, $\Gamma(x+1)=x\Gamma(x)$ and $\ln\Gamma$ convex, so that the
factorials have one natural interpolation and Euler found it. Volume~4, Chapters~5 and~9, proves that theorem and extends $\Gamma$ to the complex plane; in Volume~5 it multiplies
the zeta function.

\begin{figure}[ht]
\centering
\fig{ch06-gammafn}
\caption{Left: the Gamma function on $]0,4.5]$, with the points $(n+1,n!)$ marked. The minimum is at $x\approx1.4616$, where $\Gamma\approx0.8856$; at $x=\frac12$ the value is $\sqrt\pi\approx1.772$. Right: the integrands $t^{x-1}e^{-t}$ for $x=\frac12,1,2,3,4$; the area under each is $\Gamma(x)$. As $x$ grows the bump moves right and swells like $(x-1)!$.}
\label{fig:gammafn}
\end{figure}

\section{Numerical integration, and the next term of Stirling}

Most integrals have no closed form, and a computer evaluates them by a Riemann-type sum. Proposition~\ref{prop:inttools}(d) says the error is $O\left(\frac1n\right)$
with $n$ rectangles. Two small changes do much better: evaluate at the midpoint of each interval, or replace each rectangle by a trapezoid.

\begin{proposition}[midpoint and trapezoid rules]\label{prop:quadrature}
Let $f$ be of class $C^2$ on $[a,b]$ with $|f''|\le M$, let $h=\frac{b-a}n$ and $x_k=a+kh$. Then
\[
\begin{gathered}
\left|\int_a^bf-h\sum_{k=1}^nf\!\left(x_k-\tfrac h2\right)\right|\le\frac{M(b-a)^3}{24n^2},\\
\left|\int_a^bf-h\sum_{k=1}^n\frac{f(x_{k-1})+f(x_k)}2\right|\le\frac{M(b-a)^3}{12n^2}.
\end{gathered}
\]
\end{proposition}
\begin{proof}
It suffices to bound the error on one interval $[c-\frac h2,c+\frac h2]$ of length $h$ by $\frac{Mh^3}{24}$ (midpoint) and $\frac{Mh^3}{12}$ (trapezoid), then add the $n$ errors.

Midpoint. By Taylor--Lagrange at $c$ to order $1$, $f(x)=f(c)+f'(c)(x-c)+r(x)$ with $|r(x)|\le\frac M2(x-c)^2$. Integrating over the interval, the linear term integrates to $0$ by symmetry, so
$\left|\int f-hf(c)\right|=\left|\int r\right|\le\frac M2\int_{-h/2}^{h/2}u^2\dd u=\frac{Mh^3}{24}$.

Trapezoid. Let the interval be $[\alpha,\beta]$ with $\beta-\alpha=h$ and $E=\int_\alpha^\beta f-h\frac{f(\alpha)+f(\beta)}2$. Integrating by parts twice with the polynomial $p(x)=\frac{(x-\alpha)(\beta-x)}2$, which vanishes at both ends and has $p''=-1$:
$\int_\alpha^\beta f=\int_\alpha^\beta f\cdot(-p'')=-\left[fp'\right]_\alpha^\beta+\int_\alpha^\beta f'p'=h\,\frac{f(\alpha)+f(\beta)}2+\left[f'p\right]_\alpha^\beta-\int_\alpha^\beta f''p$, using $p'(\alpha)=\frac h2$, $p'(\beta)=-\frac h2$ and $p(\alpha)=p(\beta)=0$. So $E=-\int_\alpha^\beta f''p$. As $p\ge0$ with $\int_\alpha^\beta p=\frac{h^3}{12}$, $|E|\le M\frac{h^3}{12}$.
\end{proof}

The midpoint error bound is half the trapezoid error bound, and the two errors have opposite signs for convex $f$ (the chord lies above the curve, and the midpoint rectangle, which has the same area as the trapezoid under the tangent at the midpoint, lies below it); the combination $\frac{2M+T}3$, Simpson's rule, has error $O\left(\frac1{n^4}\right)$ if $f$ is of class $C^4$. This higher-order estimate requires more smoothness than the proposition above. A numerical calculation measures all three on $\int_0^1e^{-x^2}\dd x$ (Figure~\ref{fig:quadrature}).

\begin{figure}[ht]
\centering
\fig{ch06-quadrature}
\caption{Left: the trapezoid rule (grey) and the midpoint rule (yellow) on one interval for a convex function: the trapezoid overestimates, the midpoint rectangle, which has the same area as the tangent trapezoid, underestimates. Right: the errors of the rectangle, trapezoid, midpoint and Simpson rules for $\int_0^1e^{-x^2}\dd x$, against $n$, on logarithmic scales. The slopes are $-1$, $-2$, $-2$ and $-4$: each rule's error is a power of $\frac1n$, the proposition above supplies the midpoint and trapezoid bounds. The Simpson rate uses four continuous derivatives, which this integrand has.}
\label{fig:quadrature}
\end{figure}

\begin{proposition}[Stirling, one term further]\label{prop:stirlingnext}
\[n!=\sqrt{2\pi n}\left(\frac ne\right)^n\left(1+\frac1{12n}+O\!\left(\frac1{n^2}\right)\right).\]
\end{proposition}
\begin{proof}
Let $d_n=\ln n!-\left(n+\frac12\right)\ln n+n$ as in Exercise~\ref{ex:stirling}, with limit $d=\ln\sqrt{2\pi}$ by Theorem~\ref{thm:stirlingC}. Expanding $\ln(1+x)=x-\frac{x^2}2+\frac{x^3}3-\frac{x^4}4+O(x^5)$ at $x=\frac1n$ one order further than in that exercise,
\[
d_n-d_{n+1}=\left(n+\tfrac12\right)\ln\left(1+\tfrac1n\right)-1=\frac1{12n^2}-\frac1{12n^3}+O\!\left(\frac1{n^4}\right).
\]
Summing from $n$ to infinity, $d_n-d=\sum_{k\ge n}(d_k-d_{k+1})=\frac1{12}\sum_{k\ge n}\frac1{k^2}+O\!\left(\sum_{k\ge n}\frac1{k^3}\right)$. Now $\frac1{k(k+1)}\le\frac1{k^2}\le\frac1{(k-1)k}$ and both bounds telescope: $\frac1n\le\sum_{k\ge n}\frac1{k^2}\le\frac1{n-1}=\frac1n+O\left(\frac1{n^2}\right)$; likewise $\sum_{k\ge n}\frac1{k^3}\le\frac1{n}\sum_{k\ge n}\frac1{k(k-1)}=\frac1{n(n-1)}=O\left(\frac1{n^2}\right)$. So $d_n-d=\frac1{12n}+O\left(\frac1{n^2}\right)$, and $\frac{n!}{\sqrt{2\pi n}(n/e)^n}=e^{d_n-d}=1+\frac1{12n}+O\left(\frac1{n^2}\right)$, since $e^u=1+u+O(u^2)$.
\end{proof}

For $n=10$ the correction $1+\frac1{120}$ turns $\sqrt{2\pi}=2.5066$ into $2.5275$, and the box above measured $2.5276$. As an advanced outlook, the full asymptotic expansion continues $1+\frac1{12n}+\frac1{288n^2}-\frac{139}{51840n^3}-\cdots$; it diverges for every $n$ (the coefficients eventually grow like factorials), yet its first few terms give $n!$ to extraordinary accuracy: an \emph{asymptotic} series, a notion Volume~5 needs for the prime number theorem. The full expansion and its divergence are not proved here; see the references on Stirling's formula below.

\portraits{leibniz}{Gottfried Wilhelm Leibniz (1646--1716). Portrait by Christoph Bernhard Francke, c.\,1700. Public domain, via Wikimedia Commons.}{riemann}{Bernhard Riemann (1826--1866). Photograph, 1863. Public domain, via Wikimedia Commons.}

\section{Workshop: approximations with a safety net}
Taylor's formula has two parts: a polynomial you can calculate and a remainder you must control. Keeping the first and discarding the second is like buying a parachute and leaving the straps at home.

\begin{examplex}[how many digits of $e$?]
For $0\le x\le1$, Taylor--Lagrange at $0$ gives
\[
\left|e^x-\sum_{k=0}^{n}\frac{x^k}{k!}\right|\le\frac{e}{(n+1)!}.
\]
One may use $e<3$ without assuming the desired decimal: $k!\ge2^{k-1}$ for $k\ge2$, with strict inequality for $k\ge3$, so the exponential series gives $e<1+1+\sum_{k=2}^{\infty}2^{-(k-1)}=3$. At $n=9$ the error is below $3/10!<10^{-6}$. This is an absolute error statement, not a promise that every rounding boundary behaves conveniently.
\end{examplex}

\begin{figure}[ht]\centering\fig{review06-taylor}
\caption{Solid: $e^x$; dashed: $T_2(x)=1+x+x^2/2$, on $[0,1]$. Dotted: the upper envelope $T_2(x)+3x^3/6$. Taylor's integral remainder places the exponential between the two polynomial bounds throughout this interval.}\end{figure}

\subsection*{The integral that refuses to be elementary}
\textbf{Mission.} Approximate $I=\int_0^1e^{-x^2}\dd x$ to within $10^{-4}$ using rational numbers only. An elementary antiderivative is not available; that is no reason to go hungry.

\textbf{Hint.} Expand $e^{-u}$ at $u=0$, then put $u=x^2$. Use the integral remainder to justify the integrated error.

\textbf{Solution.} Taylor's formula gives
\[
e^{-x^2}=\sum_{k=0}^m\frac{(-1)^kx^{2k}}{k!}+r_m(x),
\qquad |r_m(x)|\le\frac{x^{2m+2}}{(m+1)!}
\]
for $0\le x\le1$, since every derivative of $e^{-u}$ has absolute value at most $1$ there. Integrating the \emph{finite} sum,
\[
\left|I-\sum_{k=0}^m\frac{(-1)^k}{k!(2k+1)}\right|
\le\frac1{(m+1)!(2m+3)}.
\]
At $m=6$ the bound is $1/75600<10^{-4}$. There is no need to claim that every function equals its Taylor series, or to interchange an infinite sum and an integral.

\begin{reminder}[local information, global caution]
An antiderivative formula on an interval must respect singularities. The antiderivatives of $1/x$ on the disconnected set $\R\setminus\{0\}$ are $\ln|x|+C_+$ on the positive half-line and $\ln|x|+C_-$ on the negative half-line, with unrelated constants. Likewise $\int_{-1}^1(1/x)\dd x$ does not converge as an improper integral: cancellation of symmetric cutoffs defines a different notion, a principal value.
\end{reminder}

\photograph{rocket}{Falcon 9 launching DSCOVR, with bright exhaust against a dark sky.}{DSCOVR launch on Falcon 9. SpaceX Photos, public domain; via scikit-image. A launch trajectory needs both estimates and error control. This photograph motivates that distinction; the one-variable examples here are not a flight-dynamics model.}


\section*{What this chapter establishes}
Differentiation and integration are inverse operations, with an explicit remainder when a function is replaced by a polynomial. Wallis's one recurrence produces the product for $\pi/2$, the constant in Stirling's formula, and the Gaussian integral. A numerical method is a theorem only when its error is bounded.

Volume~4 is the workshop that uses these formulas every day: special functions, asymptotics and sharper remainders. Volume~6 constructs the integral that this chapter took as a contract. Volume~7 extends the derivative and the transform. Volume~11 needs the Gaussian integral as the law of a random increment.

\section{Exercises}

\begin{exercise}[$\star$]\label{ex:xnlnx}
Compute $\displaystyle\int_0^1x^n\ln x\dd x$ for $n\in\N$.
\end{exercise}

\begin{exercise}[$\star$]\label{ex:riemann}
Compute the limits of $\displaystyle\sum_{k=1}^n\frac n{n^2+k^2}$ and of $\displaystyle\left(\prod_{k=1}^n\left(1+\frac kn\right)\right)^{1/n}$.
\end{exercise}

\begin{exercise}[$\star\star$]\label{ex:taylorln}
\begin{enumerate}[label=(\alph*)]
\item Show that $e^x=\sum_{k=0}^\infty\frac{x^k}{k!}$ for every real $x$.
\item Integrating $\frac1{1+t}=\sum_{k=0}^{n-1}(-t)^k+\frac{(-t)^n}{1+t}$ over $[0,1]$, give a second proof that $\sum_{k\ge1}\frac{(-1)^{k+1}}k=\ln2$, with an explicit error bound.
\item For $z=x+iy$ with $x,y\in\R$, show that
$\sum_{n\ge0}z^n/n!=e^x(\cos y+i\sin y)$, agreeing with the complex exponential
of Chapter~\ref{ch:complex}. Use the Cauchy product; no complex differentiation is needed.
\end{enumerate}
\end{exercise}

\begin{exercise}[$\star\star$]\label{ex:eirr}
Show that $e$ is irrational.
\end{exercise}

\begin{exercise}[$\star\star$]\label{ex:wallisprod}
\emph{Wallis' product.} Show that $\displaystyle\frac{W_{2p}}{W_{2p+1}}\to1$, and deduce
\[
\begin{aligned}\frac\pi2&=\lim_{p\to\infty}\prod_{k=1}^p\frac{(2k)^2}{(2k-1)(2k+1)}\eqbr=\frac21\cdot\frac23\cdot\frac43\cdot\frac45\cdots\end{aligned}
\]
\end{exercise}

\begin{exercise}[$\star\star$]\label{ex:rl}
\emph{Riemann--Lebesgue.} Let $f$ be of class $C^1$ on $[a,b]$. Show that $\displaystyle\int_a^bf(t)\sin(\lambda t)\dd t\to0$ as $\lambda\to\infty$.
\end{exercise}

\begin{exercise}[$\star\star\star$]\label{ex:gaussian}
\emph{The Gaussian integral, without polar coordinates.}
\begin{enumerate}[label=(\alph*)]
\item Show that $\left(1-\frac{x^2}n\right)^n\le e^{-x^2}$ for $0\le x\le\sqrt n$, and $e^{-x^2}\le\left(1+\frac{x^2}n\right)^{-n}$ for $x\ge0$.
\item Show that $\int_0^{\sqrt n}\left(1-\frac{x^2}n\right)^n\dd x=\sqrt n\,W_{2n+1}$ and $\int_0^\infty\left(1+\frac{x^2}n\right)^{-n}\dd x=\sqrt n\,W_{2n-2}$ (for $n\ge1$).
\item Conclude that $\displaystyle\int_0^\infty e^{-x^2}\dd x=\frac{\sqrt\pi}2$.
\end{enumerate}
\end{exercise}

\subsection*{Further exercises}

\begin{exercise}[$\star$]\label{ex:xnex}
Show that $\displaystyle\int_0^\infty x^ne^{-x}\dd x=n!$ by induction, and that $\displaystyle\int_0^1x^m(\ln x)^n\dd x=\frac{(-1)^nn!}{(m+1)^{n+1}}$ for $m,n\in\N$ (this generalises Exercise~\ref{ex:xnlnx}).
\end{exercise}

\begin{exercise}[$\star$]\label{ex:leibnizpi}
\emph{Leibniz's series (1674).} Integrating $\frac1{1+t^2}=\sum_{k=0}^{n-1}(-1)^kt^{2k}+\frac{(-1)^nt^{2n}}{1+t^2}$ over $[0,x]$, show that $\arctan x=\sum_{k=0}^\infty\frac{(-1)^kx^{2k+1}}{2k+1}$ for $0\le x\le1$, with error at most $\frac{x^{2n+1}}{2n+1}$ after $n$ terms. Deduce $\frac\pi4=1-\frac13+\frac15-\frac17+\cdots$, and explain why this is a terrible way to compute $\pi$ and why $\frac\pi4=4\arctan\frac15-\arctan\frac1{239}$ (Machin, 1706) is a good one.
\end{exercise}

\begin{exercise}[$\star\star$]\label{ex:lhopital}
\begin{enumerate}[label=(\alph*)]
\item Compute $\lim_{x\to0}\frac{x-\sin x}{x^3}$, $\lim_{x\to0}\frac{e^x-1-x}{x^2}$ each by l'H\^opital and by Taylor--Young. Compute also
$\lim_{x\to0^+}x\ln x$ by l'H\^opital and by an elementary exponential bound. Deduce $\lim_{x\to0^+}x^x$.
\item \emph{The $\frac\infty\infty$ form.} Let $f,g$ be differentiable on $]a,a+h[$ with $g'\ne0$, $g(x)\to+\infty$ as $x\to a^+$ and $\frac{f'}{g'}\to\ell\in\R$. Show that $\frac fg\to\ell$. \emph{(Fix $x_0$ close to $a$; for $x<x_0$, Cauchy's mean value theorem on $[x,x_0]$ gives $\frac{f(x)-f(x_0)}{g(x)-g(x_0)}$ close to $\ell$; then let $x\to a$ with $x_0$ fixed.)}
\end{enumerate}
\end{exercise}

\begin{exercise}[$\star\star\star$]\label{ex:trapezoid}
\emph{Bridge: Taylor approximation and roots-of-unity cancellation.}
Apply the trapezoid rule with $n=10$ to $\displaystyle\int_0^1\frac{\dd x}{1+x^2}=\frac\pi4$. Compute the error bound of Proposition~\ref{prop:quadrature} and compare with the actual error ($T_{10}\approx0.784981$, $\frac\pi4\approx0.785398$). Then show that for the \emph{periodic} integral $\int_0^{2\pi}e^{\cos x}\dd x$ the trapezoid rule with $n=10$ is already correct to eight decimals, and find which step of the error estimate becomes an identity for a periodic integrand.
\end{exercise}

\begin{exercise}[$\star\star$]\label{ex:frullani}
\emph{Frullani's integral (1828).} Let $f$ be continuous on $[0,\infty[$ with a limit $f(\infty)$ at infinity, and $a,b>0$. Show that
\[
\int_0^\infty\frac{f(ax)-f(bx)}x\dd x=\bigl(f(0)-f(\infty)\bigr)\ln\frac ba.
\]
Deduce $\int_0^\infty\frac{e^{-ax}-e^{-bx}}x\dd x=\ln\frac ba$ and $\int_0^\infty\frac{\arctan(bx)-\arctan(ax)}{x}\dd x=\frac\pi2\ln\frac ba$.
\end{exercise}

\begin{exercise}[$\star\star$]\label{ex:gammahalf}
\begin{enumerate}[label=(\alph*)]
\item Show that $\displaystyle\int_0^1(1-t^2)^n\dd t=W_{2n+1}$ and $\displaystyle\int_0^1t^{m}(1-t)^{n}\dd t=\frac{m!\,n!}{(m+n+1)!}$ for $m,n\in\N$ (Euler's Beta integral).
\item Deduce that $\displaystyle\int_0^\infty x^{2n}e^{-x^2}\dd x=\frac{(2n)!}{4^nn!}\cdot\frac{\sqrt\pi}2$. These half-line integrals determine the even moments of $e^{-x^2}/\sqrt\pi$, which Volume~3 needs.
\end{enumerate}
\end{exercise}

\begin{exercise}[$\star\star\star$]\label{ex:niven}
\emph{Niven's proof that $\pi$ is irrational (1947).} Suppose $\pi=\frac ab$ with $a,b\in\N^*$, and for $n\ge1$ let $f(x)=\frac{x^n(a-bx)^n}{n!}$ and $F(x)=f(x)-f''(x)+f^{(4)}(x)-\dots+(-1)^nf^{(2n)}(x)$.
\begin{enumerate}[label=(\alph*)]
\item Show that $f^{(k)}(0)$ and $f^{(k)}(\pi)$ are integers for every $k$.
\item Show that $\frac{\dd}{\dd x}\bigl(F'(x)\sin x-F(x)\cos x\bigr)=f(x)\sin x$, and deduce that $I_n=\int_0^\pi f(x)\sin x\dd x$ is an integer.
\item Show that $0<I_n<\frac{\pi(\pi a)^n}{n!}$, and conclude.
\end{enumerate}
\end{exercise}

\begin{exercise}[$\star\star\star$]\label{ex:sophomore}
\emph{The sophomore's dream (Johann Bernoulli, 1697).} Show that
\[
\int_0^1x^{-x}\dd x=\sum_{n=1}^\infty\frac1{n^n}=1.29128\ldots\qquad\text{and}\qquad\int_0^1x^{x}\dd x=\sum_{n=1}^\infty\frac{(-1)^{n+1}}{n^n}.
\]
\emph{(Write $x^{-x}=e^{-x\ln x}$, expand the exponential with the explicit remainder of Exercise~\ref{ex:taylorln}(a), and use Exercise~\ref{ex:xnex}.)}
\end{exercise}

\begin{challenge}
\begin{enumerate}[label=(\arabic*)]
\item \emph{The Basel problem by a double integral.} Show that $\zeta(2)=\int_0^1\!\!\int_0^1\frac{\dd x\dd y}{1-xy}$, and evaluate the double integral by the change of variables $x=\frac{u-v}{\sqrt2}$, $y=\frac{u+v}{\sqrt2}$ (rotate the square by $45^\circ$): the integrand becomes $\frac2{2-u^2+v^2}$, whose integral in $v$ is an $\arctan$. This is Apostol's 1983 proof, with Volume~6 supplying the two theorems on double integrals it needs.
\item \emph{$\pi$ is transcendental.} Lindemann proved in 1882 that $\pi$ is not a root of any polynomial with integer coefficients, which settled the ancient question of squaring the circle. His proof refines Hermite's 1873 proof for $e$, which refines the integration-by-parts trick of Exercises~\ref{ex:eirr} and~\ref{ex:niven}. Read Hermite's proof for $e$ (in Hardy and Wright, Chapter~XI) and locate the trick.
\item \emph{Open.} Nobody knows whether $e+\pi$, $e\pi$ or $\pi^e$ is irrational. Show that if $e+\pi$ and $e\pi$ were both rational, then $e$ and $\pi$ would be roots of one quadratic with rational coefficients, hence algebraic. Either conclusion contradicts a transcendence theorem: Hermite proved in 1873 that $e$ is transcendental, and Lindemann proved in 1882 that $\pi$ is. Those proofs are outlined in the notes; they are not proved in this chapter. The parenthetical quadratic step is the part to write out.
\end{enumerate}
Notes at the end of the chapter.
\end{challenge}

\begin{beauty}
One recurrence, $W_n=\frac{n-1}nW_{n-2}$, gives Wallis' product, the constant in Stirling's formula, and the Gaussian integral. These three
are really the same fact, and in Volume~3 they meet again as the central limit theorem: the $\sqrt{2\pi}$ in the normal density is Stirling's
$\sqrt{2\pi}$. The Gamma function, found by Euler to interpolate $n!$, turns out to be the function that multiplies $\zeta(s)$ in Riemann's
functional equation (Volume~5): the integers' factorials and the primes' distribution are tied together by one integral. Riemann--Lebesgue
(Exercise~\ref{ex:rl}) is the first step of Fourier analysis in Volume~4, and Niven's proof (Exercise~\ref{ex:niven}) is a classic short road
to a theorem that took humanity two thousand years to reach.
\end{beauty}

\begin{exercise}[$\star\star$]\label{pb:lightsout}
Lamps numbered by $\Z/n\Z$, $n\ge3$, stand on a circle, each on or off. Pressing lamp $i$ switches the state of the lamps
$i-1$, $i$ and $i+1$.
\begin{enumerate}[label=(\alph*)]
\item Show that every configuration can be switched off if and only if $3\nmid n$.
\item Let $3\mid n$. For $r\in\{0,1,2\}$, let $S_r$ be the number of lit lamps $i$ with $i\equiv r\pmod3$. Show that a
configuration can be switched off if and only if $S_0$, $S_1$ and $S_2$ have the same parity. What proportion of the
configurations is that?
\end{enumerate}
\end{exercise}

\begin{exercise}[$\star\star$]\label{pb:hyperplanes}
\emph{Covering a cube minus a point.}
\begin{enumerate}[label=(\alph*)]
\item (Alon and Füredi, 1993.) Show that affine hyperplanes of $\mathbb F_q^n$ that cover every point except $0$ must number at
least $n(q-1)$, and that $n(q-1)$ is enough.
\item (Jamison 1977, Brouwer and Schrijver 1978.) A \emph{blocking set} is a set of points of $\mathbb F_q^n$ that meets every affine
hyperplane. Show that a blocking set has at least $n(q-1)+1$ points. Find one of that size.
\end{enumerate}
\end{exercise}

\begin{exercise}[$\star\star$]\label{pb:kakeya}
\emph{Dvir's Kakeya theorem.} A \emph{Kakeya set} $K\subset\mathbb F_q^n$ contains a line in every direction: for every
$v\ne0$ there is $a$ with $\{a+tv:t\in\mathbb F_q\}\subset K$.
\begin{enumerate}[label=(\alph*)]
\item (Dvir, 2008.) Show that $|K|\ge\binom{q+n-1}n$. In particular $|K|\ge\frac{q^n}{n!}$: a Kakeya set fills a fixed fraction
of the space.
\item For $q$ odd, build a Kakeya set in $\mathbb F_q^2$ with $\frac{q(q+1)}2+\frac{q-1}2$ points.
\end{enumerate}
\end{exercise}

\section{Hints}

\begin{hint}{ex:xnlnx}
Integrate by parts on $[\varepsilon,1]$, differentiating $\ln x$.
\end{hint}
\begin{hint}{ex:riemann}
$\frac n{n^2+k^2}=\frac1n\cdot\frac1{1+(k/n)^2}$. For the product, take logarithms.
\end{hint}
\begin{hint}{ex:taylorln}
(a) Taylor--Lagrange with $\sup|\exp^{(n+1)}|\le e^{|x|}$ on $[0,x]$. (b) The error term is $\int_0^1\frac{t^n}{1+t}\dd t$; bound $\frac1{1+t}$ by $1$. (c) Apply Taylor to sine and cosine, split the series at $iy$
into even and odd terms, and prove the series addition law with the Cauchy product.
\end{hint}
\begin{hint}{ex:eirr}
If $e=\frac pq$, show that $q!\left(e-\sum_{k=0}^q\frac1{k!}\right)$ is an integer strictly between $0$ and $1$.
\end{hint}
\begin{hint}{ex:wallisprod}
$1\le\frac{W_{2p}}{W_{2p+1}}\le\frac{W_{2p-1}}{W_{2p+1}}$, and the last ratio is given by the recurrence.
\end{hint}
\begin{hint}{ex:rl}
Integrate by parts: $\sin(\lambda t)$ has the small antiderivative $-\frac{\cos(\lambda t)}\lambda$.
\end{hint}
\begin{hint}{ex:gaussian}
(a) $\ln(1+u)\le u$. (b) Substitute $x=\sqrt n\sin\theta$, then $x=\sqrt n\tan\theta$. (c) Squeeze, and use $W_n\sim\sqrt{\frac\pi{2n}}$.
\end{hint}
\begin{hint}{ex:xnex}
Parts, differentiating the power; for the second integral, parts differentiating $(\ln x)^n$, on $[\varepsilon,1]$.
\end{hint}
\begin{hint}{ex:leibnizpi}
The remainder is $\int_0^x\frac{t^{2n}}{1+t^2}\dd t\le\int_0^xt^{2n}\dd t$. For Machin, compute $\tan$ of $4\arctan\frac15$ with the addition formula, twice.
\end{hint}
\begin{hint}{ex:lhopital}
(a) For $x^x$, write $e^{x\ln x}$. (b) Write $\frac{f(x)}{g(x)}=\frac{f(x)-f(x_0)}{g(x)-g(x_0)}\cdot\frac{g(x)-g(x_0)}{g(x)}+\frac{f(x_0)}{g(x)}$, and note that the middle factor and the last term tend to $1$ and $0$ as $x\to a^+$.
\end{hint}
\begin{hint}{ex:trapezoid}
$f''(x)=\frac{6x^2-2}{(1+x^2)^3}$, bounded by $2$ on $[0,1]$. For the periodic case, identify the vanishing endpoint difference in the
kernel proof. To certify the stated precision, expand $e^u$ through degree $19$
at $u=\cos x$ and use the ten-point roots-of-unity filter on each power of cosine.
\end{hint}
\begin{hint}{ex:frullani}
On $[\varepsilon,A]$, substitute in each integral separately: $\int_\varepsilon^A\frac{f(ax)}x\dd x=\int_{a\varepsilon}^{aA}\frac{f(u)}u\dd u$. The difference reduces to $\int_{a\varepsilon}^{b\varepsilon}\frac{f(u)}u\dd u-\int_{aA}^{bA}\frac{f(u)}u\dd u$, where $f$ is nearly constant.
\end{hint}
\begin{hint}{ex:gammahalf}
(a) $t=\sin\theta$ for the first; for the second, integrate by parts to get a recurrence in $n$, or substitute $t=\sin^2\theta$ and use Wallis. (b) $x^2=t$ turns the integral into $\frac12\Gamma\left(n+\frac12\right)$.
\end{hint}
\begin{hint}{ex:niven}
(a) $f$ has integer coefficients divided by $n!$ and only powers $x^n,\dots,x^{2n}$; and $f(\pi-x)=f(x)$. (b) Differentiate, and note $F+F''=f$ because $f^{(2n+2)}=0$. (c) On $]0,\pi[$, $0<x(a-bx)<\pi a$ and $0<\sin x\le1$; then $\frac{c^n}{n!}\to0$.
\end{hint}
\begin{hint}{ex:sophomore}
$|x\ln x|\le\frac1e$ on $[0,1]$, so the Taylor--Lagrange remainder of the exponential after $N$ terms is at most $\frac{e^{1/e}}{(N+1)!\,e^{N+1}}$, uniformly in $x$; integrate term by term with $\int_0^1\frac{(-x\ln x)^n}{n!}\dd x=\frac1{(n+1)^{n+1}}$.
\end{hint}

\begin{hint}{pb:lightsout}
Identify configurations with $\mathbb F_2[x]/(x^n-1)$: pressing the set $u$ produces $(x^{-1}+1+x)u$. For (b), the matrix is
symmetric, so its image is the orthogonal of its kernel.
\end{hint}

\begin{hint}{pb:hyperplanes}
(a) Compare $\prod_i(1-a_i\cdot x)$ with $\prod_j\prod_{s\ne0}(x_j-s)$. (b) The hyperplanes $\{x:a\cdot x=1\}$ avoid $0$; swap the
roles of $a$ and $x$.
\end{hint}

\begin{hint}{pb:kakeya}
(a) A non-zero polynomial of degree at most $q-1$ vanishing on $K$ exists if $K$ is small. Restrict it to a line and look at the
top coefficient. (b) Use the lines $y=mx+m^2$.
\end{hint}

\section{Solutions}
\begin{solution}{ex:xnlnx}
\idea A product of a power and a logarithm: integrate by parts so as to \emph{differentiate} the logarithm, which turns it into $\frac1x$ and leaves a pure power. The integral is improper at $0$, so work on $[\varepsilon,1]$ and let $\varepsilon\to0$.
\writeup On $[\varepsilon,1]$, integrate by parts with $u'=x^n$, $v=\ln x$:
$\int_\varepsilon^1x^n\ln x\dd x=\left[\frac{x^{n+1}}{n+1}\ln x\right]_\varepsilon^1-\int_\varepsilon^1\frac{x^n}{n+1}\dd x$.
As $\varepsilon\to0$, $\varepsilon^{n+1}\ln\varepsilon\to0$, and the integral equals $-\dfrac1{(n+1)^2}$.
\end{solution}

\begin{solution}{ex:riemann}
\idea A sum of $n$ terms whose general term is a function of $\frac kn$, times $\frac1n$, is a Riemann sum. Factor out $n^2$ in the first to see $\frac1n\cdot\frac1{1+(k/n)^2}$. For a product, take logarithms to get a sum.
\writeup
\[
\sum_{k=1}^n\frac n{n^2+k^2}=\frac1n\sum_{k=1}^n\frac1{1+(k/n)^2}\to\int_0^1\frac{\mathrm dx}{1+x^2}=\frac\pi4
\]
 by Proposition~\ref{prop:inttools}(d).
For the product, take logarithms:
\[
\frac1n\sum_{k=1}^n\ln\left(1+\frac kn\right)\to\int_0^1\ln(1+x)\dd x=\left[(1+x)\ln(1+x)-x\right]_0^1=2\ln2-1,
\]
so the limit is $e^{2\ln2-1}=\frac4e$ by continuity of $\exp$.
\end{solution}

\begin{solution}{ex:taylorln}
\idea (a) Showing that a function equals its Taylor series means showing the remainder tends to $0$; for $\exp$ all derivatives are $\exp$, so Taylor--Lagrange gives a bound $\frac{|x|^{n+1}}{(n+1)!}e^{|x|}$, and factorials beat powers. (This reconciles the real exponential of analysis, the solution of $y'=y$, $y(0)=1$, with the series representation of $e^{i\theta}$, identified in part~(c) below.) (b) The identity is a finite geometric sum, so it is exact; integrating it gives the partial sum of the alternating series plus an explicit error integral, easy to bound.
\writeup (a) Apply Theorem~\ref{thm:taylor}(b) at $a=0$. Every derivative of $\exp$ is $\exp$, and $\exp\le e^{|x|}$ between $0$ and $x$, so the error is at most $\frac{|x|^{n+1}}{(n+1)!}e^{|x|}$, which tends to $0$ because the ratio of consecutive terms, $\frac{|x|}{n+2}$, tends to $0$.

(b) Integrating over $[0,1]$: $\ln2=\sum_{k=0}^{n-1}\frac{(-1)^k}{k+1}+(-1)^n\int_0^1\frac{t^n}{1+t}\dd t$, and
$0\le\int_0^1\frac{t^n}{1+t}\dd t\le\int_0^1t^n\dd t=\frac1{n+1}$. So the partial sums of $\sum\frac{(-1)^{k+1}}k$ are within $\frac1{n+1}$ of $\ln2$.

(c) Put $E(z)=\sum_{n\ge0}z^n/n!$. At $z=0$ its sum is $1$; for $z\ne0$ the ratio test proves absolute convergence.
Taylor--Lagrange for sine and cosine bounds each remainder at a real $y$ by
$|y|^{N+1}/(N+1)!$, tending to zero. Their Taylor coefficients therefore give
\[
E(iy)=\sum_{k\ge0}\frac{(-1)^k y^{2k}}{(2k)!}
+i\sum_{k\ge0}\frac{(-1)^k y^{2k+1}}{(2k+1)!}
=\cos y+i\sin y.
\]
Separating the even and odd terms is justified by absolute convergence.
The complex Cauchy product and the binomial formula give
\[
E(a)E(b)=\sum_{n\ge0}\sum_{k=0}^n
\frac{a^kb^{n-k}}{k!(n-k)!}
=\sum_{n\ge0}\frac{(a+b)^n}{n!}=E(a+b).
\]
Part (a) gives $E(x)=e^x$ for real $x$. Hence
$E(x+iy)=E(x)E(iy)=e^x(\cos y+i\sin y)$, as required.
\end{solution}

\begin{solution}{ex:eirr}
\idea Irrationality proofs follow one pattern: assume $x=\frac pq$, build an integer from it, and show that integer lies strictly between $0$ and $1$. The series of $e$ converges very fast, so the tail after $\frac1{q!}$ is tiny; multiplying by $q!$ makes the head an integer and keeps the tail below $1$.
\writeup Suppose $e=\frac pq$ with $p,q\ge1$. Then $A=q!\left(e-\sum_{k=0}^q\frac1{k!}\right)$ is an integer, since $q!e=(q-1)!\,p$ and $\frac{q!}{k!}\in\N$ for $k\le q$. But
\[
\begin{aligned}0<A&=\sum_{k>q}\frac{q!}{k!}=\frac1{q+1}+\frac1{(q+1)(q+2)}+\cdots\eqbr<\sum_{j\ge1}\frac1{(q+1)^j}=\frac1q\le1.\end{aligned}
\]
No integer lies strictly between $0$ and $1$.
\end{solution}

\begin{solution}{ex:wallisprod}
\idea The ratio $\frac{W_{2p}}{W_{2p+1}}$ is squeezed between $1$ and $\frac{W_{2p-1}}{W_{2p+1}}$ by monotonicity, and the recurrence computes the latter exactly. Then write the ratio with the closed forms (ii) and reorganise the factorials into a product over $k$.
\writeup By monotonicity, $W_{2p+1}\le W_{2p}\le W_{2p-1}$, so $1\le\frac{W_{2p}}{W_{2p+1}}\le\frac{W_{2p-1}}{W_{2p+1}}=\frac{2p+1}{2p}\to1$. By Proposition~\ref{prop:wallis}(ii),
\[
\begin{aligned}\frac{W_{2p}}{W_{2p+1}}&=\frac\pi2\cdot\frac{(2p)!\,(2p+1)!}{16^p\,(p!)^4}\eqbr=\frac\pi2\prod_{k=1}^p\frac{(2k-1)(2k+1)}{(2k)^2},\end{aligned}
\]
because $(2p)!=\prod_{k=1}^p(2k-1)(2k)$, $(2p+1)!=(2p)!\,(2p+1)$ and $\prod_{k=1}^p(2k)=2^pp!$. Letting $p\to\infty$ gives the product formula.
\end{solution}

\begin{solution}{ex:rl}
\idea A rapidly oscillating factor $\sin\lambda t$ has a small antiderivative, $-\frac{\cos\lambda t}\lambda$. Integrating by parts moves the derivative onto $f$ (which is why we assume $C^1$) and produces a factor $\frac1\lambda$.
\writeup Integrate by parts:
\[\begin{aligned}\int_a^bf(t)\sin(\lambda t)\dd t&=\left[-f(t)\frac{\cos\lambda t}\lambda\right]_a^b\eqbr+\frac1\lambda\int_a^bf'(t)\cos(\lambda t)\dd t.\end{aligned}\]
Both terms together are at most $\frac C\lambda$ in absolute value, where $C=|f(a)|+|f(b)|+\int_a^b|f'|$, so they tend to $0$.
\end{solution}

\begin{solution}{ex:gaussian}
\idea We know $\left(1\pm\frac{x^2}n\right)^{\pm n}\to e^{-x^2}$, and the integrals of these approximations are Wallis integrals after a trigonometric substitution ($1-\sin^2=\cos^2$ and $1+\tan^2=\frac1{\cos^2}$). So squeeze $e^{-x^2}$ between them, and let the known asymptotics $W_n\sim\sqrt{\frac\pi{2n}}$ finish. The inequalities in (a) are the concavity bound $\ln(1+u)\le u$.
\writeup (a) With $u=\frac{x^2}n$: $\ln(1-u)\le-u$ for $0\le u<1$ gives $n\ln(1-u)\le-x^2$, the first inequality (it is trivial at $u=1$); $\ln(1+u)\le u$ gives $(1+u)^n\le e^{nu}=e^{x^2}$, the second.

(b) Substitute $x=\sqrt n\sin\theta$, $\theta\in\left[0,\frac\pi2\right]$: $\left(1-\frac{x^2}n\right)^n\dd x=\cos^{2n}\theta\cdot\sqrt n\cos\theta\dd\theta$, giving $\sqrt nW_{2n+1}$.
Substitute $x=\sqrt n\tan\theta$, $\theta\in\left[0,\frac\pi2\right[$: $\left(1+\frac{x^2}n\right)^{-n}\dd x=\cos^{2n}\theta\cdot\frac{\sqrt n}{\cos^2\theta}\dd\theta$, giving $\sqrt nW_{2n-2}$ (an improper integral; apply the change of variables on $[0,A]$ and let $A\to\infty$).

(c) Let $I=\int_0^\infty e^{-x^2}\dd x$ (it converges since $e^{-x^2}\le e^{-x}$ for $x\ge1$). By (a) and (b),
$\sqrt nW_{2n+1}\le\int_0^{\sqrt n}e^{-x^2}\dd x\le I\le\sqrt nW_{2n-2}$. By Proposition~\ref{prop:wallis}(iii),
$\sqrt nW_{2n+1}\sim\sqrt n\sqrt{\frac\pi{4n}}=\frac{\sqrt\pi}2$, and likewise for $\sqrt nW_{2n-2}$. Hence $I=\frac{\sqrt\pi}2$.
\end{solution}


\begin{solution}{ex:xnex}
\idea Both are the integration-by-parts machine: each step lowers the exponent of the power (first integral) or of the logarithm (second) by one, so the answers are products of the successive exponents, that is, factorials. The improper ends ($\infty$ and $0$) are harmless because $x^ne^{-x}$ and $x^{m+1}(\ln x)^n$ tend to $0$ there.
\writeup Let $I_n=\int_0^\infty x^ne^{-x}\dd x$. Then $I_0=1$, and for $n\ge1$, on $[0,A]$, with $u'=e^{-x}$, $v=x^n$: $\int_0^Ax^ne^{-x}\dd x=\left[-x^ne^{-x}\right]_0^A+n\int_0^Ax^{n-1}e^{-x}\dd x$; as $A\to\infty$, $A^ne^{-A}\to0$, so $I_n=nI_{n-1}$ and $I_n=n!$.
Let $J_{m,n}=\int_0^1x^m(\ln x)^n\dd x$. Then $J_{m,0}=\frac1{m+1}$, and for $n\ge1$, on $[\varepsilon,1]$ with $u'=x^m$, $v=(\ln x)^n$: $\int_\varepsilon^1x^m(\ln x)^n\dd x=\left[\frac{x^{m+1}}{m+1}(\ln x)^n\right]_\varepsilon^1-\frac n{m+1}\int_\varepsilon^1x^m(\ln x)^{n-1}\dd x$. The bracket tends to $0$ (as $\varepsilon^{m+1}(\ln\varepsilon)^n\to0$), so $J_{m,n}=-\frac n{m+1}J_{m,n-1}$, and $J_{m,n}=\frac{(-1)^nn!}{(m+1)^{n+1}}$.
\end{solution}

\begin{solution}{ex:leibnizpi}
\idea Same recipe as for $\ln2$: a finite geometric sum with an explicit remainder, integrated. The remainder is tiny for small $x$ and merely small for $x=1$, which is why Leibniz's series is slow and Machin's fast: $\frac15$ and $\frac1{239}$ are small.
\writeup Integrating the identity from $0$ to $x$ and using $\arctan'=\frac1{1+t^2}$ (fundamental theorem): $\arctan x=\sum_{k=0}^{n-1}\frac{(-1)^kx^{2k+1}}{2k+1}+(-1)^nR_n(x)$ with $0\le R_n(x)=\int_0^x\frac{t^{2n}}{1+t^2}\dd t\le\int_0^xt^{2n}\dd t=\frac{x^{2n+1}}{2n+1}$. For $0\le x\le1$ this tends to $0$, which proves the series and the error bound. At $x=1$: $\frac\pi4=1-\frac13+\frac15-\cdots$, James Gregory's and Leibniz's series (1671, 1674), known to Madhava in Kerala around 1400.
For $\pi/4$ the error after $n$ terms is at most $1/(2n+1)$; for $\pi$ the bound is $4/(2n+1)$. To guarantee absolute error below $10^{-10}$ for $\pi$ by this bound requires $n\ge2\cdot10^{10}$. With Machin: $\tan\left(2\arctan\frac15\right)=\frac{2/5}{1-1/25}=\frac5{12}$, $\tan\left(4\arctan\frac15\right)=\frac{10/12}{1-25/144}=\frac{120}{119}$, and $\tan\left(4\arctan\frac15-\frac\pi4\right)=\frac{120/119-1}{1+120/119}=\frac1{239}$; since $4\arctan\frac15-\frac\pi4\in\,]0,\frac\pi2[$, this gives Machin's formula. The series for $\arctan\frac15$ has error at most $\frac{5^{-(2n+1)}}{2n+1}$ after $n$ terms: $7$ terms of it and $2$ of $\arctan\frac1{239}$ give a combined error in $\pi$ below $10^{-10}$ after the factors $16$ and $4$ are included. Machin computed $100$ decimals this way in 1706, and formulas of the same kind held the record for more than 270 years: the first computer calculation, $2037$ decimals on the ENIAC in 1949, used Machin's own formula, and only in 1982 (Tamura and Kanada, with Gauss's arithmetic-geometric mean) did the record pass to other methods.
\end{solution}

\begin{solution}{ex:lhopital}
\idea (a) is routine with either tool; the point is to see that Taylor--Young does the first two limits in two lines where l'H\^opital needs several differentiations each. (b) The $\frac\infty\infty$ form cannot use the trick of extending by $0$ at $a$; instead one compares $f$ and $g$ to their values at a fixed nearby point $x_0$, where Cauchy's theorem controls the increments, and then lets $g(x)\to\infty$ wash out the fixed quantities.
\writeup (a) By Taylor--Young at $0$: $x-\sin x=\frac{x^3}6+o(x^3)$, so the first limit is $\frac16$; $e^x-1-x=\frac{x^2}2+o(x^2)$, so the second is $\frac12$. By l'H\^opital, the first needs three differentiations ($\frac{1-\cos x}{3x^2}$, $\frac{\sin x}{6x}$, $\frac{\cos x}6$) and the second two. For $x\ln x=\frac{\ln x}{1/x}$, the $\frac\infty\infty$ form of (b) applies (with $x\to0^+$, $g=\frac1x\to+\infty$): $\frac{1/x}{-1/x^2}=-x\to0$. For an independent bound, put $t=-\ln x\to\infty$. Taylor's theorem gives
$e^t\ge t^2/2$ for $t>0$, so $|x\ln x|=te^{-t}\le2/t\to0$.
Hence $x^x=e^{x\ln x}\to e^0=1$.

(b) Let $\varepsilon>0$ and choose $x_0\in\,]a,a+h[$ with $\left|\frac{f'(t)}{g'(t)}-\ell\right|\le\varepsilon$ for $a<t<x_0$. For $a<x<x_0$, Cauchy's mean value theorem on $[x,x_0]$ gives $t\in\,]x,x_0[$ with $\frac{f(x)-f(x_0)}{g(x)-g(x_0)}=\frac{f'(t)}{g'(t)}$ (note $g(x)\ne g(x_0)$ by Rolle, as $g'\ne0$), so this quotient is within $\varepsilon$ of $\ell$. Now
\[
\frac{f(x)}{g(x)}=\frac{f(x)-f(x_0)}{g(x)-g(x_0)}\cdot\left(1-\frac{g(x_0)}{g(x)}\right)+\frac{f(x_0)}{g(x)}.
\]
As $x\to a^+$, $g(x)\to\infty$, so the factor in parentheses tends to $1$ and the last term to $0$, while the first quotient stays within $\varepsilon$ of $\ell$. Hence $\left|\frac{f(x)}{g(x)}-\ell\right|\le2\varepsilon$ for $x$ close enough to $a$.
\end{solution}

\begin{solution}{ex:trapezoid}
\idea The error bound is a worst case; the actual error is smaller because $f''$ changes sign on $[0,1]$. For a periodic integrand, extracting the mean of the error kernel produces a leading term proportional to $f'(b)-f'(a)=0$. For a certified ten-point estimate, a different language is more efficient:
Taylor approximation plus the roots-of-unity filter makes the error nearly vanish
term by term. This combines complex algebra with quantitative calculus.
\writeup With $f(x)=\frac1{1+x^2}$: $f'(x)=\frac{-2x}{(1+x^2)^2}$ and $f''(x)=\frac{6x^2-2}{(1+x^2)^3}$, which lies between $-2$ (at $0$) and $\frac12$ (at $1$), so $M=2$. The bound is $\frac{2\cdot1}{12\cdot100}=\frac1{600}\approx0.00167$; the actual error is $\frac\pi4-T_{10}\approx0.000417$, four times smaller, because $f''$ is negative near $0$ and positive near $1$, and the two parts of the error partly cancel.

For the periodic case, first identify the cancellation in the error kernel.
On each subinterval write $p=h^2/12+\widetilde p$, where
$\int\widetilde p=0$. The constant parts of $-\int f''p$ sum to
$-h^2(f'(2\pi)-f'(0))/12=0$ for $f(x)=e^{\cos x}$.
This explains an improvement, but does not yet certify eight decimals.
Roots of unity give a sharper bound.

\emph{How to find the certificate.} The ten equally spaced points annihilate
oscillations $e^{imx}$ unless $10$ divides $m$. Taylor turns $e^{\cos x}$ into
a finite sum of such oscillations, so only a few terms can contribute to the error.

\emph{Write-up.} Let $P(u)=\sum_{r=0}^{19}u^r/r!$. For $-1\le u\le1$,
Taylor--Lagrange and $e<3$ give $|e^u-P(u)|<3/20!$.
Write $I(g)=\int_0^{2\pi}g(x)\,\dd x$ and
$T(g)=(2\pi/10)\sum_{k=0}^9g(2\pi k/10)$; periodicity makes this the
trapezoid rule. The approximation changes $T(g)-I(g)$ by less than $12\pi/20!$.
By the binomial formula,
\[
\cos^r x=2^{-r}\sum_{j=0}^r\binom rj e^{i(r-2j)x}.
\]
The integral of $e^{imx}$ is zero for nonzero integer $m$, by integrating
its real and imaginary parts. The ten-point average is zero unless $10\mid m$,
by the roots-of-unity filter. Since $r\le19$, only frequencies $10$ and $-10$
remain in the difference, giving
\[
T(P\circ\cos)-I(P\circ\cos)
=4\pi\sum_{j=0}^4\frac1{2^{10+2j}j!(10+j)!}.
\]
Now $(10+j)!\ge10!\,11^j$, hence this positive sum is at most
\[
\frac{4\pi}{2^{10}10!}\sum_{j=0}^4\frac1{44^j j!}
<\frac{16}{2^{10}10!}\frac{44}{43},
\]
using $\pi<4$ and a geometric-series bound. Here $\pi<4$ follows already from
$\pi/4=\int_0^1(1+x^2)^{-1}\,\dd x<1$.
Together with the uniform remainder,
\[
|T(e^{\cos x})-I(e^{\cos x})|
<\frac{16\cdot44}{43\cdot2^{10}10!}+\frac{48}{20!}
<4.42\times10^{-9}<\frac12\,10^{-8}.
\]
This is the promised eight-decimal absolute error guarantee. Numerically,
$T_{10}=7.954926524\ldots$ and $I=7.954926521\ldots$; the certificate, not the
agreement of two decimals, establishes the accuracy. No infinite sum has been
interchanged with an integral.
\end{solution}

\begin{solution}{ex:frullani}
\idea The two integrals $\int\frac{f(ax)}x$ and $\int\frac{f(bx)}x$ diverge separately (at $0$ if $f(0)\ne0$, at $\infty$ if $f(\infty)\ne0$), but after the substitution $u=ax$ and $u=bx$ they become the same integral $\int\frac{f(u)}u\dd u$ over slightly different ranges, and the difference is two short integrals, one near $0$ and one near $\infty$, where $f$ is nearly constant and $\int\frac{\dd u}u=\ln\frac ba$.
\writeup For $0<\varepsilon<A$, substituting $u=ax$ and $u=bx$ (Proposition~\ref{prop:inttools}(c)),
\[
\int_\varepsilon^A\frac{f(ax)-f(bx)}x\dd x=\int_{a\varepsilon}^{aA}\frac{f(u)}u\dd u-\int_{b\varepsilon}^{bA}\frac{f(u)}u\dd u=\int_{a\varepsilon}^{b\varepsilon}\frac{f(u)}u\dd u-\int_{aA}^{bA}\frac{f(u)}u\dd u
\]
(assume $a<b$; the two integrals over $[b\varepsilon,aA]$ cancel). On $[a\varepsilon,b\varepsilon]$, $f(u)=f(0)+o(1)$ uniformly as $\varepsilon\to0$ (continuity at $0$), so the first integral is $\bigl(f(0)+o(1)\bigr)\ln\frac ba$; likewise the second is $\bigl(f(\infty)+o(1)\bigr)\ln\frac ba$ as $A\to\infty$. Letting $\varepsilon\to0$ and $A\to\infty$ gives the formula, and in particular shows that the improper integral converges.
With $f(x)=e^{-x}$: $f(0)-f(\infty)=1$, giving $\ln\frac ba$. With $f(x)=-\arctan x$: $f(0)-f(\infty)=\frac\pi2$, giving $\frac\pi2\ln\frac ba$.
\end{solution}

\begin{solution}{ex:gammahalf}
\idea These integrals all reduce to Wallis integrals by the substitutions $t=\sin\theta$ or $t=\sin^2\theta$, which turn $1-t^2$ or $1-t$ into a $\cos^2$. The Beta integral is best found by a recurrence in $n$ obtained by parts, exactly as for $W_n$. The Gaussian moments come from $\Gamma$ by $x^2=t$.
\writeup (a) With $t=\sin\theta$, $\dd t=\cos\theta\dd\theta$: $\int_0^1(1-t^2)^n\dd t=\int_0^{\pi/2}\cos^{2n+1}\theta\dd\theta=W_{2n+1}$.
Let $B(m,n)=\int_0^1t^m(1-t)^n\dd t$. Then $B(m,0)=\frac1{m+1}$, and for $n\ge1$, integrating by parts with $u'=t^m$, $v=(1-t)^n$: $B(m,n)=\left[\frac{t^{m+1}}{m+1}(1-t)^n\right]_0^1+\frac n{m+1}B(m+1,n-1)=\frac n{m+1}B(m+1,n-1)$. Iterating, $B(m,n)=\frac{n}{m+1}\cdot\frac{n-1}{m+2}\cdots\frac1{m+n}\cdot B(m+n,0)=\frac{n!\,m!}{(m+n)!}\cdot\frac1{m+n+1}=\frac{m!\,n!}{(m+n+1)!}$.

(b) With $t=x^2$, $\dd t=2x\dd x$: $\int_0^\infty x^{2n}e^{-x^2}\dd x=\frac12\int_0^\infty t^{n-1/2}e^{-t}\dd t=\frac12\Gamma\left(n+\frac12\right)=\frac{(2n)!}{4^nn!}\cdot\frac{\sqrt\pi}2$ by Proposition~\ref{prop:gammafn}(b). For $n=0$ this is the Gaussian integral; for $n=1$, $\frac{\sqrt\pi}4$. The density $e^{-x^2}/\sqrt\pi$ on $\R$ has total mass $1$, so its even moments are
\[
\frac{2}{\sqrt\pi}\int_0^\infty x^{2n}e^{-x^2}\dd x=\frac{(2n)!}{4^nn!}.
\]
The product $1\cdot3\cdot5\cdots(2n-1)=\frac{(2n)!}{2^nn!}$ is the even moment of the standard normal density $(2\pi)^{-1/2}e^{-x^2/2}$. It follows from the integral above by the substitution $x=u/\sqrt2$, not by dividing by the mass of $e^{-x^2}$.
\end{solution}

\begin{solution}{ex:niven}
\idea The pattern of Exercise~\ref{ex:eirr} once more: from the assumption $\pi=\frac ab$ build an integer, here $I_n$, and show it lies strictly between $0$ and $1$ for $n$ large. The polynomial $f$ is designed so that its derivatives at $0$ and at $\pi$ are integers (the $\frac1{n!}$ is exactly cancelled by the $k!$ coming out of differentiation), and the alternating sum $F$ is designed so that $F''+F=f$, which makes $\int f\sin x$ computable by parts in one stroke.
\writeup (a) Expand $f(x)=\frac1{n!}\sum_{j=n}^{2n}c_jx^j$ with integers $c_j$ (the coefficients of $x^n(a-bx)^n$). Then $f^{(k)}(0)=\frac{k!\,c_k}{n!}$ for $n\le k\le2n$ and $0$ otherwise; since $k\ge n$, $\frac{k!}{n!}$ is an integer, and so is $f^{(k)}(0)$. Next, $f(\pi-x)=\frac{(\pi-x)^n(a-b\pi+bx)^n}{n!}=\frac{(\pi-x)^n(bx)^n}{n!}=f(x)$ because $a=b\pi$; differentiating $k$ times, $f^{(k)}(\pi-x)=(-1)^kf^{(k)}(x)$, so $f^{(k)}(\pi)=(-1)^kf^{(k)}(0)$ is an integer.

(b) $\frac{\dd}{\dd x}\bigl(F'\sin x-F\cos x\bigr)=F''\sin x+F'\cos x-F'\cos x+F\sin x=(F''+F)\sin x$, and $F''+F=f$ because the series defining $F$ telescopes under $F''$: $F''=f''-f^{(4)}+\dots+(-1)^nf^{(2n+2)}$ and $f^{(2n+2)}=0$ ($\deg f=2n$). Hence
$I_n=\bigl[F'(x)\sin x-F(x)\cos x\bigr]_0^\pi=F(\pi)+F(0)$, an integer by (a).

(c) On $]0,\pi[$, $x>0$, $a-bx=b(\pi-x)>0$ and $\sin x>0$, so $f(x)\sin x>0$ and $I_n>0$. Also $x(a-bx)\le\pi a$ (as $x\le\pi$ and $a-bx\le a$) and $\sin x\le1$, so $I_n\le\pi\frac{(\pi a)^n}{n!}$. Since $\frac{c^n}{n!}\to0$ for every $c$ (ratio test), $I_n<1$ for $n$ large: an integer strictly between $0$ and $1$. So $\pi$ is irrational.
Niven's note in the \emph{Bulletin of the AMS} of 1947 is one page long; Lambert had proved the theorem in 1768 with continued fractions, in thirty.
\end{solution}

\begin{solution}{ex:sophomore}
\idea Expand $x^{-x}=e^{-x\ln x}$ as a series in powers of $x\ln x$ and integrate term by term; each term is the integral of Exercise~\ref{ex:xnex}, and the sum collapses to $\sum\frac1{n^n}$. The interchange of sum and integral is justified exactly as in Exercise~\ref{ex:eulerbasel}: the remainder of the exponential series is bounded uniformly on $[0,1]$, because $|x\ln x|\le\frac1e$ there (the maximum of $-x\ln x$ is at $x=\frac1e$).
\writeup Put $u=-x\ln x\in\left[0,\frac1e\right]$ on $]0,1]$ (extend continuously by $0$ at $x=0$). Taylor--Lagrange for $\exp$ at $0$, $e^u=\sum_{k=0}^N\frac{u^k}{k!}+r_N(u)$ with $0\le r_N(u)\le\frac{u^{N+1}}{(N+1)!}e^{u}\le\frac{e^{1/e}}{(N+1)!\,e^{N+1}}=:\rho_N$, a bound independent of $x$. Integrating over $[0,1]$,
\[
\int_0^1x^{-x}\dd x=\sum_{k=0}^N\frac1{k!}\int_0^1(-x\ln x)^k\dd x+\theta_N\rho_N,\qquad0\le\theta_N\le1.
\]
By Exercise~\ref{ex:xnex},
\[
\int_0^1(-x\ln x)^k\dd x=(-1)^k\int_0^1x^k(\ln x)^k\dd x=\frac{k!}{(k+1)^{k+1}}.
\]
So the sum is $\sum_{k=0}^N\frac1{(k+1)^{k+1}}$, that is $\sum_{n=1}^{N+1}\frac1{n^n}$, and letting $N\to\infty$ (with $\rho_N\to0$) gives the first formula. For the second, replace $u$ by $-u=x\ln x$: the same bound holds, and $\int_0^1(x\ln x)^k\dd x=\frac{(-1)^kk!}{(k+1)^{k+1}}$ gives $\sum_{n\ge1}\frac{(-1)^{n+1}}{n^n}=0.78343\ldots$ The first series converges so fast ($\frac1{10^{10}}$ at $n=10$) that ten terms give $1.2912859971$. Bernoulli published both in 1697; the nickname ``sophomore's dream'' suggests the dream that integrals and sums may be swapped at will, which here, for once, comes true.
\end{solution}

\subsection*{Notes on the challenges}

(1) Integrate the finite geometric identity. Its remainder is $R_N=(xy)^N/(1-xy)$, and on $[0,1]^2$ away from the corner,
$1-xy\ge\sqrt{(1-x)(1-y)}$. Indeed, $1-xy$ is at least each of $1-x$ and $1-y$.
Consequently
\[
0\le\int_0^1\!\!\int_0^1 R_N\dd x\dd y
\le\left(\int_0^1\frac{x^N}{\sqrt{1-x}}\dd x\right)^2
=4W_{2N+1}^2\longrightarrow0.
\]
The last equality uses $x=\sin^2t$. This bound justifies the limiting interchange independently of the series being evaluated, and gives
$\int_0^1\!\!\int_0^1(1-xy)^{-1}\dd x\dd y=\sum_{n\ge0}(n+1)^{-2}$.
After rotation the region is $0\le u\le\sqrt2$, $|v|\le\min(u,\sqrt2-u)$. Integrating in $v$ gives
$\frac4{\sqrt{2-u^2}}\arctan\frac{\min(u,\sqrt2-u)}{\sqrt{2-u^2}}$.
Put $u=\sqrt2\sin\theta$. The arctangent is $\theta$ up to $\theta=\pi/6$, and $\pi/4-\theta/2$ thereafter. Thus the integral is
\[
4\int_0^{\pi/6}\theta\dd\theta+
4\int_{\pi/6}^{\pi/2}\left(\frac\pi4-\frac\theta2\right)\dd\theta=\frac{\pi^2}6.
\]
Apostol's note is in \emph{Math.\ Intelligencer} 5 (1983), 59--60; Volume~6 supplies Fubini's theorem and the change-of-variables theorem used here.

(2) Hermite's proof that $e$ is transcendental (1873) begins with a hypothetical integer relation $\sum_{k=0}^n a_ke^k=0$, with $a_0\ne0$, and uses
$f(x)=x^{p-1}(x-1)^p\cdots(x-n)^p/(p-1)!$ for a large prime $p$.
Put $F=f+f'+\cdots+f^{(\deg f)}$. Since $(e^{-x}F(x))'=-e^{-x}f(x)$,
\[
I_k=\int_0^k e^{k-x}f(x)\dd x=e^kF(0)-F(k).
\]
The individual $I_k$ need not be integers. The relation makes their weighted sum an integer:
$\sum_{k=0}^n a_kI_k=-\sum_{k=0}^n a_kF(k)$.
For a prime $p>n$ not dividing $a_0$, the derivative values are integers, $F(k)\equiv0\pmod p$ for $1\le k\le n$, and
$F(0)\equiv(-1)^{np}(n!)^p\not\equiv0\pmod p$.
Thus this integer is nonzero. On the fixed interval $[0,n]$, the numerator of $f$ is bounded by $C^p$ for a constant $C$ independent of $p$; division by $(p-1)!$ makes the weighted integral tend to zero, a contradiction.
Lindemann's extension to $\pi$ uses $e^{i\pi}=-1$ and algebraic numbers in place of $0,\dots,n$. Hardy and Wright, Chapter~XI, gives both proofs. The transcendence of $\pi$ settles squaring the circle, since every constructible number is algebraic (Volume~2).

(3) If $e+\pi=r$ and $e\pi=s$ were rational, $e$ and $\pi$ would be the roots of $X^2-rX+s$, hence algebraic, contradicting Hermite and Lindemann. The irrationality of each of $e+\pi$, $e\pi$, $\pi^e$, $\pi^{\sqrt2}$ remains open; by contrast $e^\pi$ is transcendental (Gelfond 1929), and so is $e^{\pi\sqrt{163}}=262537412640768743.99999999999925\ldots$, whose nearness to an integer Volume~12 explains.

\begin{solution}{pb:lightsout}
\idea Pressing commutes with rotating the circle, so the pressing map is a circulant matrix over $\mathbb F_2$, that is,
multiplication by a polynomial in the ring $\mathbb F_2[x]/(x^n-1)$. It is invertible exactly when that polynomial is prime
to $x^n-1$. For (b), a symmetric matrix has image orthogonal to its kernel, and the kernel is easy to guess.

\writeup Encode a set $u$ of lamps as $u(x)=\sum_{i\in u}x^i$ in $R=\mathbb F_2[x]/(x^n-1)$. Pressing the lamps of $u$
switches each lamp as many times as it has pressed lamps among itself and its neighbours, so it adds
$(x^{-1}+1+x)u=x^{-1}(1+x+x^2)u$ to the configuration (order does not matter, and pressing a lamp twice does nothing). As
$x$ is a unit in $R$, we must understand multiplication by $f=1+x+x^2$.

\emph{Lemma.} Over a field, multiplication by $f$ on $K[x]/(x^n-1)$ has a kernel of dimension $\deg g$, with
$g=\gcd(f,x^n-1)$. Indeed, $uf\equiv0$ means $x^n-1\mid uf$, that is $h\mid u\frac fg$ with $h=\frac{x^n-1}g$, and since
$\gcd(h,\frac fg)=1$ this means $h\mid u$. The multiples of $h$ of degree $<n$ are $hv$ with $\deg v<\deg g$.

(a) $1+x+x^2$ has no root in $\mathbb F_2$, so it is irreducible, and $g$ is $1$ or $1+x+x^2$. Since
$x^3-1=(x-1)(x^2+x+1)$, $x^3\equiv1$ modulo $f$, while $x\not\equiv1$ and $x^2\not\equiv1$. So $f\mid x^n-1$ if and only if
$3\mid n$. When $3\nmid n$, the kernel is $0$, pressing is a bijection of $R$, and every configuration can be switched off.
When $3\mid n$, the kernel has dimension $2$, so some configurations cannot.

(b) Let $3\mid n$, and let $v_r$ be the set of lamps $i\not\equiv r\pmod3$. Pressing $v_r$ switches each lamp twice or not at
all: a lamp in $v_r$ has exactly one neighbour in $v_r$, and a lamp outside it has both neighbours in it. So $v_0,v_1,v_2$
are in the kernel, which they span (two of them are independent and the kernel has dimension $2$). The pressing matrix
$A$ is symmetric, and over any field the image of $A$ is the orthogonal of the kernel of $A^T$ for the dot product (a
dimension count, Volume~0, Chapter~7, which holds over $\mathbb F_2$ as well). So a configuration $c$ can be switched off if
and only if $c\cdot v_r=0$ for $r=0,1,2$, that is $S_1+S_2$, $S_0+S_2$, $S_0+S_1$ are all even, that is $S_0\equiv S_1\equiv S_2
\pmod2$. The image has dimension $n-2$, so this is $\frac14$ of the $2^n$ configurations.
\tools Volume~2, Chapters~6 and~9 (polynomial rings, finite fields); Volume~0, Chapter~7 (rank and orthogonality).
\leads The commercial game \emph{Lights Out} (1995) is the $5\times5$ grid version. Its pressing matrix has a kernel of
dimension $2$ too, so only a quarter of the boards are solvable. Sutner (1989) showed that on every finite graph pressing
all vertices is enough to reach the all-off state from all-on: a theorem about $\mathbb F_2$-linear algebra.
\end{solution}

\begin{solution}{pb:hyperplanes}
\idea If fewer than $n(q-1)$ hyperplanes did the job, the product of their equations would agree, on all of $\mathbb F_q^n$, with
a polynomial whose leading monomial is $x_1^{q-1}\cdots x_n^{q-1}$, and the Nullstellensatz forbids this.

\writeup (a) The $n(q-1)$ hyperplanes $x_j=s$, $1\le j\le n$, $s\in\mathbb F_q^*$, cover every $x\ne0$ (some $x_j\ne0$) and miss $0$.
Suppose $m<n(q-1)$ hyperplanes $H_i=\{x:a_i\cdot x=b_i\}$ cover $\mathbb F_q^n\setminus\{0\}$ and miss $0$; then $b_i\ne0$ and we
may take $b_i=1$. Let
\[
f(x)=\prod_{i=1}^m(1-a_i\cdot x)-c\prod_{j=1}^n\prod_{s\in\mathbb F_q^*}(x_j-s),
\]
with $c$ chosen so that $f(0)=0$ (the second product is non-zero at $0$). For $x\ne0$, $x$ lies on some $H_i$, so the first product
vanishes, and some $x_j\in\mathbb F_q^*$, so the second does too. Thus $f$ vanishes on $\mathbb F_q^n$. But $\deg f=n(q-1)$, since
$m<n(q-1)$, and the coefficient of $x_1^{q-1}\cdots x_n^{q-1}$ is $-c\ne0$. Problem~\ref{pb:nullstellensatz}(b) with $S_j=\mathbb F_q$ and
$t_j=q-1$ gives a point where $f\ne0$: a contradiction.

(b) Let $B$ be a blocking set; translating, we may assume $0\in B$. Every hyperplane $\{x:a\cdot x=1\}$, $a\ne0$, misses $0$, so it
contains some $y\in B\setminus\{0\}$. In other words the sets $\{a:a\cdot y=1\}$, for $y\in B\setminus\{0\}$, are hyperplanes (in the
space of the $a$'s) missing $0$ and covering every $a\ne0$. By (a), $|B|-1\ge n(q-1)$. The union of the $n$ coordinate axes, with
$1+n(q-1)$ points, is blocking: a hyperplane $a\cdot x=b$ with $b\ne0$ meets the axis $j$ for any $a_j\ne0$, and one with $b=0$
contains $0$.

\tools Problem~\ref{pb:nullstellensatz}; Volume~0, Chapter~7 (hyperplanes).
\leads Over $\R$ the cube $\{0,1\}^n$ minus a point needs $n$ hyperplanes, and the same proof works. Projective blocking sets are
different: in a projective plane of order $q$ a line has $q+1$ points, and the real question is how small a blocking set
containing no line can be (Blokhuis 1994: at least $\frac32(q+1)$ for $q$ prime).
\end{solution}

\begin{solution}{pb:kakeya}
\idea Linear algebra gives a low-degree polynomial vanishing on a small $K$. On each line of $K$ it vanishes at $q$ points, more
than its degree, so it vanishes identically there, and its top coefficient along the line, which is the top homogeneous part
evaluated at the direction, is $0$. A polynomial of low degree that vanishes at every direction is $0$.

\writeup (a) The polynomials of degree at most $q-1$ in $n$ variables form a space of dimension $\binom{q-1+n}n$, the number of
monomials. If $|K|<\binom{q+n-1}n$, the linear map $f\mapsto(f(x))_{x\in K}$ has a non-trivial kernel: some $f\ne0$ of degree
$d\le q-1$ vanishes on $K$. Since $K\ne\emptyset$, $f$ is not a non-zero constant, so $d\ge1$. Let $f_d\ne0$ be its homogeneous part of
degree $d$. For $v\ne0$, take $a$ with $a+tv\in K$ for all $t$. Then $t\mapsto f(a+tv)$ is a polynomial in $t$ of degree at most $d<q$
with $q$ roots, so it is $0$, and its coefficient of $t^d$, which is $f_d(v)$, is $0$. Also $f_d(0)=0$. So $f_d$ vanishes on
$\mathbb F_q^n$ while $\deg_{x_i}f_d\le q-1<q$; by Problem~\ref{pb:nullstellensatz}(a), $f_d=0$. Contradiction. Finally
$\binom{q+n-1}n=\frac{q(q+1)\cdots(q+n-1)}{n!}\ge\frac{q^n}{n!}$.

(b) For $m\in\mathbb F_q$ let $L_m$ be the line $y=mx+m^2$, of slope $m$, and let $K$ be the union of the $L_m$ and the vertical line
$x=0$. A point $(x,y)$ lies on some $L_m$ when $m^2+mx-y=0$ has a root, that is when $x^2+4y$ is a square (zero included). For fixed
$x$, $y\mapsto x^2+4y$ is a bijection of $\mathbb F_q$, and $\frac{q+1}2$ elements are squares; so $\bigcup L_m$ has $\frac{q(q+1)}2$
points. On $x=0$ these are the points $(0,y)$ with $y$ a square, $\frac{q+1}2$ of them, so the vertical line adds $\frac{q-1}2$.

\tools Volume~0, Chapter~7 (dimension counting); Problem~\ref{pb:nullstellensatz}(a). Part~(a) is also Volume~2, Exercise~M.10; the construction of (b) is new here.
\leads Dvir's proof (2008) took one page and settled a problem that Wolff had posed in 1999 as a model for the real Kakeya
conjecture, still open in dimensions $\ge4$ (Wang and Zahl announced dimension $3$ in 2025). For $n=2$ and $q$ odd the
construction (b) is optimal (Blokhuis and Mazzocca, 2008); for $q=3$ numerical evaluation confirms by brute force that the minimum is $7$,
one more than Dvir's bound $6$.
\end{solution}

\section*{Sources and further reading}

\begin{itemize}\raggedright
\item M.~Spivak, \emph{Calculus}, 4th ed., Publish or Perish, 2008, Chapters~13--20. Integrals, Taylor's theorem and the irrationality of $e$ and $\pi$, with superb exercises.
\item E.~Hairer and G.~Wanner, \emph{Analysis by Its History}, Springer, 1996, Chapter~II. Wallis, Taylor, Lagrange's remainder, the Bernoulli--l'H\^opital contract and Euler's Gamma function, with the original texts.
\item J.~Wallis, \emph{Arithmetica Infinitorum}, Oxford, 1656; English translation with commentary by J.~A.~Stedall, \emph{The Arithmetic of Infinitesimals}, Springer, 2004. The original infinite product, found by interpolation, and the quarrel with Fermat in Stedall's introduction.
\item B.~Taylor, \emph{Methodus Incrementorum Directa et Inversa}, London, 1715, Proposition~VII; J.-L.~Lagrange, \emph{Th\'eorie des fonctions analytiques}, Paris, 1797, where the remainder appears (Part~I, \S~34--36 in the 1813 edition).
\item A.~de Moivre, \emph{Approximatio ad summam terminorum binomii $(a+b)^n$ in seriem expansi}, privately printed, 12 November 1733; reprinted in the 1738 edition of \emph{The Doctrine of Chances}. The first appearance of $e^{-x^2}$. P.-S.~Laplace, \emph{M\'emoire sur la probabilit\'e des causes par les \'ev\'enements}, 1774, for the value of the Gaussian integral.
\item C.~Truesdell, \emph{The New Bernoulli Edition}, Isis 49 (1958), 54--62, and the introduction to \emph{Der Briefwechsel von Johann Bernoulli}, vol.~1, Birkh\"auser, 1955. The l'H\^opital--Bernoulli contract of 1694 and its discovery in 1921.
\item P.~J.~Davis, \emph{Leonhard Euler's integral: a historical profile of the Gamma function}, Amer.\ Math.\ Monthly 66 (1959), 849--869. Euler's letter to Goldbach of 13 October 1729, Bohr--Mollerup, and everything else about $\Gamma$.
\item I.~Niven, \emph{A simple proof that $\pi$ is irrational}, Bull.\ Amer.\ Math.\ Soc.\ 53 (1947), 509. Exercise~\ref{ex:niven}, as it was printed. J.-B.~de Stainville, \emph{M\'elanges d'analyse alg\'ebrique et de g\'eom\'etrie}, Paris, 1815, pp.~339--341, reports Fourier's proof that $e$ is irrational (Exercise~\ref{ex:eirr}).
\item G.~H.~Hardy and E.~M.~Wright, \emph{An Introduction to the Theory of Numbers}, 6th ed., Oxford University Press, 2008, Chapter~XI. The irrationality and transcendence of $e$ and $\pi$ (Hermite, Lindemann).
\item J.~Borwein, D.~Bailey and R.~Girgensohn, \emph{Experimentation in Mathematics}, A~K Peters, 2004, Section~1.5. The sophomore's dream and Bernoulli's 1697 paper (\emph{Acta Eruditorum}, May 1697).
\item K.~Conrad, \emph{Stirling's formula} and \emph{The Gaussian integral}, expository notes on the author's web page (\url{https://kconrad.math.uconn.edu/blurbs/}). Several proofs of each, to compare with Theorem~\ref{thm:stirlingC} and Exercise~\ref{ex:gaussian}.
\item P.~Lax and M.~Terrell, \emph{Calculus with Applications}, 2nd ed., Springer, 2014. A modern calculus text with Wallis, Stirling and many numerical illustrations, including the trapezoid and Simpson rules of Proposition~\ref{prop:quadrature}.
\end{itemize}
